% Chapter 13 of The Replica Rosetta Stone: counting critical points (Kac--Rice and the spiked tensor model).
\newcommand{\kacriceCrt}{\mathrm{Crt}}
\newcommand{\kacriceHess}{\operatorname{Hess}}
\newcommand{\kacricegrad}{\operatorname{grad}}
\newcommand{\kacriceGS}{\mathrm{GS}}
\newcommand{\kacriceML}{\lambda_{\mathrm{ML}}}
\newcommand{\kacriceGOE}{\mathrm{GOE}}
\newcommand{\kacriceinf}{E_\infty}

\rstchapter{Counting Critical Points}{Kac--Rice and the spiked tensor model}{TAP complexity\dict Kac--Rice and the second moment}{ch:kacrice}
\stonefoot{partial}
\begin{refsection}[chapters/13-kacrice.bib]

\begin{chaptersummary}
Maximum likelihood for tensor PCA minimizes a random polynomial on a sphere whose noise part is the pure spherical $p$-spin glass. Physicists count its critical points as Thouless--Anderson--Palmer states: they average the Kac--Rice integral, treat the determinant of the Hessian as self-averaging, and find exponentially many local minima between the ground state $-E_0$ and a threshold $-\kacriceinf$, where the minima become marginal and gradient dynamics is predicted to stop. We prove the mean counts exactly at every size, by a Gaussian regression and an identity that turns the absolute determinant of a GOE matrix into the density of one more eigenvalue, and in the limit by large deviations of the extreme eigenvalue; that the typical count equals the mean below the threshold is Subag's second-moment theorem. Exhaustive enumeration at small $N$ matches the exact formula; gradient descent at $N\le100$ stops above $-\kacriceinf$, where the exact formula places almost all minima at these sizes; and warm-started descent on the spiked tensor lands on the overlap that the landscape predicts. As of 2026 the typical landscape near the spike and the endpoint of gradient descent are not proved.
\end{chaptersummary}

\section{The question: how rugged is the likelihood of tensor PCA?}\label{sec:kacrice:question}

In \term{tensor PCA}\index{tensor PCA} one observes a $p$-tensor, the product of an unknown vector with itself $p$ times plus Gaussian noise, and wants the vector back. Take $x^*$ on the sphere $S_N=\{\sigma\in\R^N:\norm\sigma^2=N\}$ and observe
\begin{equation}\label{eq:kacrice:data}
Y=\lambda N^{-(p-1)/2}\,(x^*)^{\otimes p}+G,
\end{equation}
where $G$ has independent $\Normal(0,1)$ entries and $\lambda\ge0$ is the signal-to-noise ratio. Since $\norm{\sigma^{\otimes p}}$ is the same for all $\sigma\in S_N$, the log-likelihood of a candidate $\sigma$ is, up to a constant, a multiple of $\ip{Y}{\sigma^{\otimes p}}$, and the \term{maximum likelihood estimator}\index{maximum likelihood estimator!tensor PCA} minimizes over $S_N$ the random function
\begin{equation}\label{eq:kacrice:spiked}
H_{N,\lambda}(\sigma)=H_N(\sigma)-\lambda N\,m(\sigma)^p,\qquad m(\sigma)=\frac{\ip{\sigma}{x^*}}N,\qquad H_N(\sigma)=\frac1{N^{(p-1)/2}}\sum_{i_1,\dots,i_p=1}^NJ_{i_1\cdots i_p}\sigma_{i_1}\cdots\sigma_{i_p},
\end{equation}
with $J=-G$. The noise part $H_N$ is the \term{pure spherical $p$-spin}\index{spherical p-spin@spherical $p$-spin model} glass in the normalization of \textcite{AuffingerBenArousCerny2013} that this book uses: it is a centered Gaussian process with $\E H_N(\sigma)H_N(\sigma')=N\,R(\sigma,\sigma')^p$, where $R(\sigma,\sigma')=\ip{\sigma}{\sigma'}/N$, and energies are reported per spin. The case $\lambda=0$ is the pure model; the overlap $m$ with the truth is the book's $m$.

For $p=2$ the estimator is the top eigenvector of a spiked matrix (\cref{ch:spiked}). For $p\ge3$ it is hard to compute but statistically strong: it correlates with $x^*$ as soon as $\lambda$ exceeds a constant \parencite{JagannathLopattoMiolane2020,PerryWeinBandeira2020}. Efficient methods need far more. Gradient flow and Langevin dynamics started at random reach the spike in time of order one when $\lambda=N^\alpha$ with $\alpha>(p-2)/2$ \parencite[Theorem~2.1]{BenArousGheissariJagannath2020}, and sum-of-squares relaxations succeed from $\lambda$ of order $N^{(p-2)/4}$ \parencite{HopkinsShiSteurer2015}. Why does descent on a function whose minimizer is good fail for a whole range of $\lambda$ growing with $N$? The physicist's answer is that the landscape is rugged: it has exponentially many local minima, and descent stops in one of them.

To make this precise we count. A point $\sigma\in S_N$ is a \term{critical point}\index{critical point} of a smooth function on $S_N$ if its gradient along the sphere vanishes; its \term{index}\index{critical point!index} is the number of negative eigenvalues of its Hessian along the sphere (\cref{sec:kacrice:proof} gives the formulas). For a Borel set $B\subset\R$ let $\kacriceCrt_N(B)$ be the number of critical points of $H_N$ with $H_N(\sigma)/N\in B$, and $\kacriceCrt_{N,k}(B)$ the number of those of index $k$; local minima have index $0$. The exponential growth rate $\frac1N\log\E\kacriceCrt_N(B)$ is the \term{annealed complexity}\index{complexity!annealed}, and $\frac1N\E\log\kacriceCrt_N(B)$ the \term{quenched complexity}\index{complexity!quenched}, which describes a typical landscape.\sidenote{Physicists call the logarithm of the number of states per degree of freedom the complexity, or configurational entropy, of a glass. The word has nothing to do with computational complexity, though this chapter is about the link between the two.}

\begin{physicsbox}[Crisanti and Sommers 1995; Kurchan and Laloux 1996; Cugliandolo and Kurchan 1993; Ros, Ben Arous, Biroli and Cammarota 2019]
Let $N\to\infty$, and let $E_\infty=2\sqrt{(p-1)/p}$ and $J(x)=\frac x4\sqrt{x^2-4}-\log\frac{x+\sqrt{x^2-4}}2$ for $x\ge2$, $J(x)=0$ for $0\le x\le2$.
\begin{enumerate}
\item Critical points of the pure model ($\lambda=0$) with energy near $u$ number $e^{N\Sigma(u)}$, both on average and typically, with
\[
\Sigma(u)=\tfrac12\log(p-1)-\frac{p-2}{4(p-1)}\,u^2-J\bigl(2\abs u/\kacriceinf\bigr).
\]
\item Below the \term{threshold energy}\index{threshold energy} $-\kacriceinf$ almost all critical points are local minima, and above it they are saddles of index proportional to $N$. Local minima fill the band $[-E_0,-\kacriceinf]$, where $-E_0$, the zero of $\Sigma$, is the ground state. The most numerous minima sit at $-\kacriceinf$ and are \term{marginal}\index{marginal stability}: their Hessian has eigenvalues down to $0$. Saddles of index $k$ appear only above an energy $-E_k$, with $E_0>E_1>\dots\to\kacriceinf$. For $p=3$, $E_0=1.65700$ and $\kacriceinf=1.63299$.
\item Gradient descent from a random start ends at the threshold $-\kacriceinf$, not at the ground state.
\item For the spiked model the band of threshold minima near the equator $m=0$ survives at every fixed $\lambda$. A local minimum correlated with $x^*$ appears at $\lambda_c(p)$ and becomes the global minimum, the estimator, at $\kacriceML(p)>\lambda_c(p)$; for $p=3$, $\lambda_c=2/\sqrt3$ and $\kacriceML\approx1.21$.
\end{enumerate}
\end{physicsbox}

Parts~(i) and~(ii) are the TAP complexity of \textcite{CrisantiSommers1995}, with the index layers of \textcite{KurchanLaloux1996}; part~(iii) is the dynamical prediction of \textcite{CugliandoloKurchan1993}; part~(iv) is the replicated Kac--Rice computation of \textcite{RosEtAl2019}, translated into this book's normalization. How do physicists derive these statements, and what is proved? This chapter proves the mean counts of~(i) and~(ii), for all critical points and for minima, exactly at every $N$ and in the limit; the typical counts are Subag's theorem, part~(iii) is a conjecture whose Langevin version \cref{ch:dmft} discusses, and part~(iv) is proved for the mean counts and for the estimator, not for the typical landscape.

\section{The physics: counting TAP states}\label{sec:kacrice:physics}

\subsection{States are critical points}\label{sec:kacrice:tap}

At low temperature a physicist describes a spin glass by its \term{states}, the local minima of a free energy of the local magnetizations, the \term{TAP free energy}\index{TAP free energy} of \textcite{ThoulessAndersonPalmer1977}. Counting them is old: \textcite{BrayMoore1980} counted the solutions of the TAP equations of the Sherrington--Kirkpatrick model with the delta-function formula of \cref{sec:toolkit:kacrice}, averaged over the couplings. For the spherical $p$-spin, \textcite{KurchanParisiVirasoro1993} and \textcite{CrisantiSommers1995} wrote the TAP free energy for magnetizations $\sqrt q\,\sigma$, with $\sigma\in S_N$ and $0\le q<1$, in a form that separates radius and direction; per spin, in our normalization and with the physicists' inverse temperature $\beta$ \parencite[eq.~(6.3)]{AuffingerBenArousCerny2013},
\begin{equation}\label{eq:kacrice:tap}
f_{\mathrm{TAP}}(q,\sigma)=\frac{q^{p/2}}{\sqrt2}\,\frac{H_N(\sigma)}N-\frac1{2\beta}\log(1-q)-\frac\beta4\bigl(1+(p-1)q^p-pq^{p-1}\bigr).
\end{equation}
The direction enters only through the energy $H_N(\sigma)/N$. A stationary point therefore has $\sigma$ a critical point of $H_N$ on the sphere, and $q$ a root of a scalar equation that involves only that energy. \Textcite[Section~6]{AuffingerBenArousCerny2013} show that each critical point of $H_N$ below a temperature-dependent energy carries exactly two TAP solutions. Counting TAP states is counting critical points of the energy, and at zero temperature the states are the local minima of $H_N$.\sidenote{In \eqref{eq:kacrice:tap} the first term is the energy of the magnetization, the second is minus the temperature times the entropy $\frac12\log(1-q)$ of the fluctuations around it, and the third is the Onsager reaction term.}

\subsection{The annealed count}\label{sec:kacrice:annealed-physics}

\begin{heuristic}[The count by a delta function]
Write the number of critical points with energy in $[u,u+\dd u]$ as $\int_{S_N}\delta(\kacricegrad H_N)\abs{\det\kacriceHess H_N}\,\delta(H_N/N-u)\dd\sigma$ and average over the couplings. The law of $H_N$ is invariant under rotations, so the integrand has the same mean at every point, and the count is the area $\abs{S_N}$ times a one-point quantity.

At a point, the covariance function $N^{1-p}\ip{\sigma}{\sigma'}^p$ and its derivatives give the joint law of the value, the tangential gradient and the tangential second derivatives: the value has variance $N$, the gradient has independent coordinates of variance $p$, the second derivatives form a matrix $G$ with $\E G_{ij}G_{kl}=\frac{p(p-1)}N(\delta_{ik}\delta_{jl}+\delta_{il}\delta_{jk})$, and the three are uncorrelated. The sphere adds the Lagrange multiplier $pH_N/N$ on the diagonal, so at energy $u$ the Hessian is
\[
\sqrt{p(p-1)}\,\bigl(A-t(u)I\bigr),\qquad t(u)=u\sqrt{\frac p{p-1}}=\frac{2u}{\kacriceinf},
\]
with $A$ a GOE matrix of size $N-1$, semicircle on $[-2,2]$. The density of the gradient at $0$ is $(2\pi p)^{-(N-1)/2}$ and that of the energy is $\sqrt{N/(2\pi)}\,e^{-Nu^2/2}$.

The physicist now treats the spectrum of $A$ as self-averaging: $\E\abs{\det(A-t)}=\E\exp\sum_i\log\abs{\lambda_i-t}\approx e^{N\Phi(t)}$, with $\Phi(t)=\int\log\abs{y-t}\,\sigma_{\mathrm{sc}}(\dd y)$ the logarithmic potential of the semicircle law. Stirling's formula turns the prefactor $\abs{S_N}(2\pi p)^{-(N-1)/2}(p(p-1))^{(N-1)/2}$ into $e^{\frac N2(1+\log(p-1))+o(N)}$, and the complexity at energy $u$ is
\begin{equation}\label{eq:kacrice:sigma}
\Sigma(u)=\tfrac12\bigl(1+\log(p-1)\bigr)-\frac{u^2}2+\Phi\bigl(t(u)\bigr).\qedhere
\end{equation}
\end{heuristic}

\Cref{comp:kacrice:potential} evaluates $\Phi$ and turns \eqref{eq:kacrice:sigma} into the formula of the physics box.

\begin{computation}{The complexity from the semicircle}\label{comp:kacrice:potential}
\step \emph{The potential inside the bulk.} The Stieltjes transform $s(z)=\int\sigma_{\mathrm{sc}}(\dd y)/(y-z)=\frac12(-z+\sqrt{z^2-4})$ (\cref{thm:toolkit:semicircle}) has real part $-t/2$ at $z=t+\mathrm i0$, $\abs t<2$, so $\Phi'(t)=-\mathrm{Re}\,s(t+\mathrm i0)=t/2$ there. With $\Phi(0)=\int\log\abs y\,\sigma_{\mathrm{sc}}(\dd y)=-\frac12$ (substitute $y=2\cos\theta$), $\Phi(t)=\frac{t^2}4-\frac12$ for $\abs t\le2$.
\step \emph{Outside.} For $t>2$, $\Phi'(t)=-s(t)=\frac12(t-\sqrt{t^2-4})$, so $\Phi(t)=\frac{t^2}4-\frac12-\frac12\int_2^t\sqrt{y^2-4}\dd y=\frac{t^2}4-\frac12-J(t)$, with $J$ the rate function of the largest eigenvalue of a GOE matrix (\cref{rem:toolkit:goeldp}); $\Phi$ is even.
\step \emph{The complexity.} Since $t(u)^2=\frac p{p-1}u^2$, \eqref{eq:kacrice:sigma} becomes $\Sigma(u)=\frac12\log(p-1)-\frac{p-2}{4(p-1)}u^2-J(\abs{t(u)})$, the formula of the box. It is largest at $u=0$, where $\Sigma(0)=\frac12\log(p-1)$.
\step \emph{The index.} At energy $u$ the Hessian spectrum fills $\sqrt{p(p-1)}\,[-2-t,2-t]$ with semicircle density, so a fraction $\sigma_{\mathrm{sc}}((-2,t))$ of its eigenvalues is negative. For $t<-2$, that is $u<-\kacriceinf$, the typical critical point is a local minimum; at $u=-\kacriceinf$ the spectrum touches $0$.
\step \emph{Numbers for $p=3$.} $\Sigma(0)=\frac12\log2=0.34657$; at the threshold $\Sigma(-\kacriceinf)=\frac12\log2-\frac13=0.013240$; and $\Sigma(u)=0$ at $u=-1.656998$, the ground state $-E_0$.
\end{computation}

Below the threshold the typical critical point is a minimum, so the complexity of minima is $\Sigma_0(u)=\Sigma(u)$ for $u\le-\kacriceinf$; above it, minima would need the whole semicircle on the positive side, and physicists conclude that there are none on the exponential scale. A saddle of index $k$ below the threshold needs $k$ eigenvalues of $A$ outside the semicircle, below $t$, each at the cost $J(\abs t)$, so its complexity is $\Sigma(u)-kJ(\abs{t(u)})$ and vanishes at an energy $-E_k$ above $-E_0$ \parencite{KurchanLaloux1996}. $\Sigma_0$ increases up to $-\kacriceinf$, so the most numerous minima sit at the threshold, where the lowest Hessian eigenvalues touch zero: they are marginal, flat in a growing number of directions. That is the landscape behind part~(iii) of the box: descent slows down among marginal minima and stops where they are typical. \Cref{fig:kacrice:complexity} shows the curves.

\begin{figure}[t]
\begin{fullwidth}
\centering
\begin{tikzpicture}
\begin{groupplot}[group style={group size=2 by 1,horizontal sep=18mm},rplot,width=82mm,height=52mm]
\nextgroupplot[xlabel={energy per spin $u$},ylabel={complexity},xmin=-1.7,xmax=0,ymin=-0.05,ymax=0.38,
  ytick={0,0.1,0.2,0.3},title={(a) all critical points, $p=3$},legend pos=south east,legend style={font=\scriptsize}]
\addplot[guide,forget plot] coordinates {(-1.6329932,-0.05) (-1.6329932,0.38)};
\addplot[guide,forget plot] coordinates {(-1.7,0) (0,0)};
\addplot[physpred] table[x=u,y=density] {figures/data/13-kacrice-complexity.dat};
\addplot[proofpred] table[x=u,y=total] {figures/data/13-kacrice-complexity.dat};
\legend{physics $\Sigma(u)$,theorem $\Theta_3(u)$}
\node[font=\scriptsize,text=muted,anchor=south west,rotate=90] at (axis cs:-1.645,0.02) {$-E_\infty$};
\nextgroupplot[xlabel={energy per spin $u$},ylabel={complexity},xmin=-1.664,xmax=-1.626,ymin=-0.008,ymax=0.017,
  xtick={-1.66,-1.65,-1.64,-1.63},ytick={-0.005,0,0.005,0.01,0.015},title={(b) the bottom: minima and saddles},
  yticklabel style={/pgf/number format/.cd,fixed,precision=3},legend pos=north west,legend style={font=\scriptsize}]
\addplot[guide,forget plot] coordinates {(-1.664,0) (-1.626,0)};
\addplot[guide,forget plot] coordinates {(-1.6329932,-0.008) (-1.6329932,0.017)};
\addplot[physpred] table[x=u,y=g] {figures/data/13-kacrice-cs95.dat};
\addplot[proofpred] table[x=u,y=minima] {figures/data/13-kacrice-complexity.dat};
\addplot[proof!60,line width=0.8pt,densely dotted,no marks] table[x=u,y=idx1] {figures/data/13-kacrice-complexity.dat};
\addplot[proof!40,line width=0.8pt,dashdotted,no marks] table[x=u,y=idx2] {figures/data/13-kacrice-complexity.dat};
\legend{minima: Crisanti--Sommers,minima: theorem,index $1$,index $2$}
\node[font=\scriptsize,text=muted,anchor=north] at (axis cs:-1.6570,-0.0005) {$-E_0$};
\node[font=\scriptsize,text=muted,anchor=north] at (axis cs:-1.6529,-0.0035) {$-E_1$};
\node[font=\scriptsize,text=muted,anchor=north] at (axis cs:-1.6503,-0.0005) {$-E_2$};
\node[font=\scriptsize,text=muted,anchor=north west] at (axis cs:-1.6330,0.0165) {$-E_\infty$};
\end{groupplot}
\end{tikzpicture}
\end{fullwidth}
\caption{The complexity of the pure spherical $3$-spin. (a) The physicists' complexity $\Sigma(u)$ of critical points at energy $u$, \eqref{eq:kacrice:sigma} (carnelian), and the limit $\Theta_3(u)$ of $\frac1N\log\E\kacriceCrt_N((-\infty,u))$ proved in \cref{thm:kacrice:annealed} (lapis dashed); they coincide on $u\le0$. (b) The bottom of the landscape: the complexity of local minima of \textcite{CrisantiSommers1995}, as printed in \textcite[eq.~(2.18)]{AuffingerBenArousCerny2013} and converted to our energies by $u=\sqrt2E$ (carnelian), the theorem for minima (lapis dashed), and the complexities of saddles of index $1$ and $2$ \parencite[Theorem~2.5]{AuffingerBenArousCerny2013}. All curves agree above the threshold $-E_\infty=-1.63299$; minima disappear below $-E_0=-1.65700$, saddles of index $k$ below $-E_k$ ($E_1=1.65287$, $E_2=1.65025$). Script: \texttt{figures/scripts/13-kacrice-complexity.py}.}\label{fig:kacrice:complexity}
\end{figure}

\subsection{The typical count and the ground state}\label{sec:kacrice:replicas}

The mean count can be dominated by rare landscapes with many more critical points than typical. Physicists compute the typical count with replicas: $\E\log\kacriceCrt=\lim_{n\to0}(\E\kacriceCrt^n-1)/n$, and $\E\kacriceCrt^n$ counts $n$-tuples of critical points, which interact through their mutual overlaps $q_{ab}$. If the dominant tuples are mutually orthogonal, $q_{ab}=0$ for $a\ne b$ (a replica-symmetric ansatz with zero overlap), then $\E\kacriceCrt^n\approx(\E\kacriceCrt)^n$ and the quenched complexity equals the annealed one; for the pure $p$-spin below the threshold this is the physicists' expectation \parencite{CrisantiSommers1995,CrisantiLeuzziRizzo2003}. In the same picture the ground state is the zero of $\Sigma_0$. A second physical derivation gives the same number: in the zero-temperature limit of the solution of \textcite{CrisantiSommers1992} with one step of replica-symmetry breaking, the ground state is $-\inf_{c,d>0}\frac12[c+pd+\frac1c\log(1+\frac cd)]$, in the form \textcite[eq.~(5.11)]{AuffingerBenArousCerny2013} derive from Talagrand's Parisi formula. For $p=3$ the infimum, at $c=0.62502$ and $d=0.34399$, is $1.6569983635$, the zero of $\Sigma_0$ to ten digits.\sidenote{Some computations drop the absolute value of the determinant. Then each critical point is counted with the sign $(-1)^{\text{index}}$; below an energy the signed count is the Euler characteristic of the sublevel set, and over all energies it is that of the sphere, $2$ or $0$ (\cref{ex:kacrice:euler}). Below the threshold, where almost all critical points are minima, the signed and unsigned counts agree on the exponential scale. \Textcite{CavagnaGiardinaParisi1998} counted stationary points of every index of the TAP free energy.}

\subsection{The spike}\label{sec:kacrice:spike-physics}

The same count works with the spike, with three changes at a point of overlap $m$. The energy is shifted by $-\lambda m^p$. The gradient is shifted by $\lambda pm^{p-1}$ times the tangential part of $x^*$, which costs the Gaussian factor $e^{-Np\lambda^2m^{2p-2}(1-m^2)/2}$ in the density of the gradient at $0$. And the Hessian acquires the rank-one term $-\sqrt{p(p-1)}\,\theta(m)ee^\top$ along the tangential direction $e$ of $x^*$, with
\[
\theta(m)=\lambda\sqrt{p(p-1)}\,m^{p-2}(1-m^2).
\]
When $\theta>1$ this term pulls an eigenvalue of $A-\theta ee^\top$ out of the semicircle, to $-(\theta+1/\theta)$: the BBP transition of \cref{ch:spiked}, which can turn a would-be minimum into a saddle. Repeating the heuristic with these changes, and adding the entropy $\frac12\log(1-m^2)$ of the points at overlap $m$, gives the complexity at overlap $m$ and energy $u$:
\begin{equation}\label{eq:kacrice:sigmalambda}
\Sigma_\lambda(m,u)=\tfrac12\bigl(1+\log(p-1)\bigr)+\tfrac12\log(1-m^2)-\tfrac p2\lambda^2m^{2p-2}(1-m^2)-\tfrac12(u+\lambda m^p)^2+\Phi\bigl(t(u)\bigr).
\end{equation}
For $p\ge3$ the spike has no gradient and no curvature at the equator $m=0$, where $\Sigma_\lambda(0,u)=\Sigma(u)$: the threshold minima of the pure model survive there for every $\lambda$. Away from the equator, $\sup_u\Sigma_\lambda(m,u)$ can touch zero again at an overlap $m^*$ close to $1$: an isolated minimum correlated with the truth. This happens exactly when $m^{2p-4}(1-m^2)=1/(p\lambda^2)$ has a root, that is for
\begin{equation}\label{eq:kacrice:lambdac}
\lambda\ge\lambda_c(p)=\Bigl(\frac{(p-1)^{p-1}}{p\,(p-2)^{p-2}}\Bigr)^{1/2},
\end{equation}
$\lambda_c(3)=2/\sqrt3=1.15470$, and $m^*$ is the largest root. \Textcite{RosEtAl2019} computed the typical complexity by a replicated Kac--Rice formula at the replica-symmetric level. For $p=3$ they find a band of minima around the equator that grows with $\lambda$; at $\lambda_c$ the isolated minimum detaches from the top of the band, at $\kacriceML\approx1.21$ it becomes the global minimum, and the band shrinks but never disappears; the most numerous minima are always the threshold minima at the equator. Their two thresholds are the inverses of the dynamical and static transition temperatures of the pure model in our normalization ($T_d=\sqrt3/2$ and $T_s=0.8288$, which are $\sqrt2$ times the values $\sqrt{3/8}$ and $0.586$ of \cref{ch:dmft}).\sidenote{\Textcite[eq.~(58)]{RosEtAl2019} state this as $r_c=\frac p2\beta_d$ and $r_{\mathrm{1ST}}=\frac p2\beta_s$ in their normalization, in which $r=p\lambda/\sqrt2$ and temperatures are those of \cref{ch:dmft}.}

The derivations above contain five uncontrolled steps: the annealed average, where the typical count is wanted; the replacement of $\E\abs{\det}$ by the exponential of the mean logarithm, which exchanges an expectation with the exponential of an extensive quantity; the continuation $n\to0$ and the zero-overlap ansatz for the typical count; and the identification of where descent stops with where minima are most numerous.

\begin{figure}[t]
\begin{fullwidth}
\centering
\begin{tikzpicture}
\begin{groupplot}[group style={group size=2 by 1,horizontal sep=18mm},rplot,width=82mm,height=52mm]
\nextgroupplot[xlabel={overlap $m$ with the truth},ylabel={complexity of minima},xmin=0,xmax=1,ymin=-0.04,ymax=0.016,
  ytick={-0.04,-0.02,0,0.01},yticklabel style={/pgf/number format/.cd,fixed,precision=2},
  title={(a) local minima of $H_{N,\lambda}$, $p=3$},legend pos=south west,legend style={font=\scriptsize}]
\addplot[guide,forget plot] coordinates {(0,0) (1,0)};
\addplot[proof!35,line width=0.9pt,no marks] table[x=m,y=min0] {figures/data/13-kacrice-spiked.dat};
\addplot[proof!60,line width=0.9pt,no marks] table[x=m,y=min1] {figures/data/13-kacrice-spiked.dat};
\addplot[proof!80,line width=0.9pt,no marks] table[x=m,y=min2] {figures/data/13-kacrice-spiked.dat};
\addplot[proof,line width=1.0pt,no marks] table[x=m,y=min3] {figures/data/13-kacrice-spiked.dat};
\legend{$\lambda=0.8$,$\lambda=\lambda_c=1.155$,$\lambda=\kacriceML=1.207$,$\lambda=2$}
\nextgroupplot[xlabel={signal-to-noise ratio $\lambda$},ylabel={final overlap $m$},xmin=0.3,xmax=6.3,ymin=-0.05,ymax=1.05,
  title={(b) gradient descent, $N=80$},legend pos=south east,legend style={font=\scriptsize}]
\addplot[guide,forget plot] coordinates {(1.1547,-0.05) (1.1547,1.05)};
\addplot[proofpred] table[x=lam,y=mstar] {figures/data/13-kacrice-mstar.dat};
\addplot[measured,ebars] table[x=lam,y=warm,y error plus=warmhalf,y error minus=warmhalf] {figures/data/13-kacrice-spikedgd.dat};
\addplot[measured3,ebars] table[x=lam,y=frac,y error plus=frachi,y error minus=fraclo] {figures/data/13-kacrice-spikedgd.dat};
\legend{$m^*(\lambda)$,warm start: mean $m$,random start: share with $m>\frac12$}
\end{groupplot}
\end{tikzpicture}
\end{fullwidth}
\caption{The spiked tensor model, $p=3$. (a) The annealed complexity of local minima at overlap $m$, maximized over the energy (\cref{thm:kacrice:bammn}), for four values of $\lambda$ (darker for larger $\lambda$). At $m=0$ every curve equals $\frac12\log2-\frac13=0.0132$, the threshold minima of the pure model. For $\lambda\ge\lambda_c$ a second branch touches zero at $m^*(\lambda)$: $0.707$ at $\lambda_c$, $0.803$ at $\kacriceML$, $0.953$ at $\lambda=2$. (b) Projected gradient descent at $N=80$ on two independent noise tensors with six starts of each kind per tensor (seed $1331$): from warm starts with $m=0.5$, the mean final overlap with a $95\%$ Student-$t$ interval over the $12$ runs; from uniform random starts, the share of runs that end with $m>\frac12$, with $95\%$ Clopper--Pearson intervals. The curve is $m^*(\lambda)$, which above $\kacriceML$ is also the overlap of the maximum likelihood estimator (\cref{thm:kacrice:jlm}); the dashed line marks $\lambda_c$. Script: \texttt{figures/scripts/13-kacrice-spiked.py}.}\label{fig:kacrice:spiked}
\end{figure}

\section{The stone}\label{sec:kacrice:stone}

\begin{stone}{the TAP count of Crisanti and Sommers, with replicas}{Auffinger, Ben Arous and \v{C}ern\'y; Subag; Ben Arous, Mei, Montanari and Nica}
\stonestep{Count}{Write the number of states as $\int\delta(\nabla H)\abs{\det\nabla^2H}\,\delta(H-Nu)$ over the sphere and average it over the couplings.}{The Kac--Rice formula (\cref{thm:toolkit:kacrice}) gives $\E\kacriceCrt_{N,k}(B)$ as an integral of a conditional expectation; its hypotheses, nondegeneracy and a modulus of continuity, make the count of nondegenerate critical points exact.}
\stonestep{The Hessian}{By isotropy the Hessian at a point is a GOE matrix shifted by $p$ times the energy, independent of the gradient.}{\Cref{lem:kacrice:regression}: at a point the value, the gradient and the second derivatives are independent Gaussians, exactly at every $N$; given the energy $u$ the Hessian is $\sqrt{p(p-1)}(A-t(u)I)$.}
\stonestep{The determinant}{The spectrum is self-averaging: $\E\abs{\det(A-t)}\approx e^{N\Phi(t)}$, with $\Phi$ the logarithmic potential of the semicircle.}{\Cref{lem:kacrice:extra}: $\E\abs{\det(A-t)}$ is $c_Ne^{Nt^2/4}$ times the density of an eigenvalue of a $\kacriceGOE(N)$ matrix at $t$, an exact identity from the joint law of the eigenvalues. No physical counterpart; the constant comes from $p=2$, where the count is $2N$ (\cref{lem:kacrice:two}).}
\stonestep{The saddle point}{Maximize the exponent over the energy window.}{The semicircle law inside the bulk and the large deviations of the extreme eigenvalue outside it (\cref{thm:kacrice:annealed}); Laplace's method on the energy is replaced by monotonicity.}
\stonestep{The threshold}{Minima need the shifted semicircle on the positive side: $u\le-\kacriceinf$; at $-\kacriceinf$ they are marginal.}{The rate of the smallest eigenvalue is infinite inside the bulk, so minima above $-\kacriceinf$ are rarer in mean than every exponential (\cref{cor:kacrice:groundstate}); ABC prove a bound $e^{-cN^2}$ on their probability.}
\stonestep{The typical count}{Replicas with zero overlap between distinct critical points: $\E\log\kacriceCrt=\log\E\kacriceCrt$.}{Subag's second moment (\cref{thm:kacrice:subag}): a two-point Kac--Rice formula shows $\E\kacriceCrt^2/(\E\kacriceCrt)^2\to1$ below $-\kacriceinf$. Two replicas suffice, and nothing is continued in $n$.}
\stonestep{The ground state}{The zero of the complexity of minima, confirmed by the zero-temperature 1RSB free energy.}{The first moment gives $\liminf\kacriceGS_N\ge-E_0$ (\cref{cor:kacrice:groundstate}); the matching bound needs the second moment (Subag) or the Parisi formula (ABC, even $p$).}
\stonestep{The spike}{Shifted energy and gradient and a rank-one Hessian term whose BBP outlier destabilizes minima; replicated Kac--Rice for the typical landscape.}{The same Gaussian regression (\cref{prop:kacrice:spiked}) and the large deviations of the outlier give the mean landscape (\cref{thm:kacrice:bammn}); the estimator through the Parisi formula (\cref{thm:kacrice:jlm}). The typical landscape near the spike has no counterpart yet.}
\end{stone}

The translation is closer here than in most chapters of the book. The physicist's integral is the Kac--Rice formula, and the covariance computation that makes the Hessian a shifted GOE matrix is the same calculation on both sides. The proof replaces two steps. The self-averaging of the determinant becomes an exact identity at every $N$, which is stronger than the physicist needs. The saddle point over the energy becomes a large-deviation estimate for one eigenvalue, the only place where the proof needs more than the semicircle law. The physicist's \enquote{no minima above the threshold} becomes a statement about the rate function of that eigenvalue: it is infinite inside the bulk.

The two routes part at the typical count. The replica computation continues in the number of replicas and assumes the dominant tuples orthogonal; the proof computes exactly two replicas, which is all that concentration needs, and the answer is the same for the pure model. For mixed models the physicist's assumption fails, and the mean count is not the typical one (\cref{sec:kacrice:fails}).

\section{The proof: the mean number of critical points}\label{sec:kacrice:proof}

We prove the mean counts of the pure model exactly at every $N$ (\cref{thm:kacrice:exact}, Theorems~2.1 and~2.2 of \textcite{AuffingerBenArousCerny2013}) and their exponential rates for all critical points and for minima (\cref{thm:kacrice:annealed}, their Theorems~2.8 and~2.5 with $k=0$), given the Kac--Rice formula of \cref{thm:toolkit:kacrice}, the facts about GOE matrices in \cref{ch:toolkit}, and the joint density of GOE eigenvalues. We follow their proof with two changes, which are the book's: we fix the normalizing constant by the case $p=2$ instead of Selberg's integral, and we treat only the cases that need the large deviations of the largest eigenvalue alone.

\paragraph{Setting.} Fix integers $p\ge2$ and $N\ge2$, and extend $H_N$ by \eqref{eq:kacrice:spiked} to a homogeneous polynomial of degree $p$ on $\R^N$, whose covariance is $N^{1-p}\ip{\sigma}{\sigma'}^p$ for all $\sigma,\sigma'$. For a smooth $F$ on $\R^N$ and $\sigma\in S_N$, the gradient of $F$ along the sphere is $P_\sigma\nabla F(\sigma)$, with $P_\sigma=I-\sigma\sigma^\top/N$, and its Hessian along the sphere is the quadratic form on $T_\sigma=\sigma^\perp$
\begin{equation}\label{eq:kacrice:hess}
\kacriceHess F(\sigma)=P_\sigma\nabla^2F(\sigma)P_\sigma-\frac{\ip{\sigma}{\nabla F(\sigma)}}N\,P_\sigma,
\end{equation}
because the second derivative of $F$ along the great circle $s\mapsto\sigma\cos(s/\sqrt N)+\sqrt N\,v\sin(s/\sqrt N)$, with $v\perp\sigma$ a unit vector, is $\ip{v}{\nabla^2F(\sigma)v}-\ip{\sigma}{\nabla F(\sigma)}/N$. For $H_N$, Euler's identity gives $\ip{\sigma}{\nabla H_N(\sigma)}=pH_N(\sigma)$. Throughout, $W\sim\kacriceGOE(N)$ is a GOE matrix as in \cref{def:toolkit:goe} (semicircle on $[-2,2]$), with ordered eigenvalues $\lambda_{(0)}\le\dots\le\lambda_{(N-1)}$ and empirical spectral measure $\mu_N$; $\rho_N$ is the density of the mean measure $\E\mu_N$. We write $A$ for an $(N-1)\times(N-1)$ symmetric matrix whose entries $A_{ab}$, $a\le b$, are independent centered Gaussians with variances $(1+\delta_{ab})/N$; that is, $A=((N-1)/N)^{1/2}W'$ with $W'\sim\kacriceGOE(N-1)$, and its density on the entries is proportional to $e^{-N\tr(A^2)/4}$. Finally $\kacriceinf=2\sqrt{(p-1)/p}$, $t(u)=2u/\kacriceinf=u\sqrt{p/(p-1)}$, and $J$ is the rate function of \cref{rem:toolkit:goeldp}.

\paragraph{The eigenvalue at the edge.} By \cref{rem:toolkit:goeldp}, $\lambda_{\max}(W)$ satisfies a large deviation principle with speed $N$ and the good rate function $J$, which is $+\infty$ on $(-\infty,2)$ and continuous and strictly increasing on $[2,\infty)$, with $J(2)=0$. Applied to closed half-lines and to open intervals, the principle gives, for $x\ge2$ and $\delta>0$,
\begin{gather}
\limsup_{N\to\infty}\frac1N\log\Prob\bigl(\lambda_{\max}(W)\ge x\bigr)\le-J(x),\qquad\liminf_{N\to\infty}\frac1N\log\Prob\bigl(x<\lambda_{\max}(W)<x+\delta\bigr)\ge-J(x),\label{eq:kacrice:ldp}\\
\lim_{N\to\infty}\frac1N\log\Prob\bigl(\lambda_{\max}(W)\le y\bigr)=-\infty\quad\text{for every }y<2.\notag
\end{gather}
Since $-W\sim\kacriceGOE(N)$, the same holds for $-\lambda_{(0)}$.

\begin{lemma}[Gaussian regression at a point]\label{lem:kacrice:regression}
Fix $\sigma\in S_N$ and an orthonormal basis $e_1,\dots,e_{N-1}$ of $T_\sigma$. The value $H_N(\sigma)$, the gradient $g=(\ip{e_i}{\nabla H_N(\sigma)})_{i<N}$ and the matrix $G=(\ip{e_i}{\nabla^2H_N(\sigma)e_j})_{i,j<N}$ are independent, with $H_N(\sigma)\sim\Normal(0,N)$, $g\sim\Normal(0,pI_{N-1})$, and $G$ distributed as $\sqrt{p(p-1)}\,A$. Consequently, in the basis $(e_i)$, $\kacriceHess H_N(\sigma)=\sqrt{p(p-1)}\,(A-t(u)I)$ with $u=H_N(\sigma)/N$ and $A$ independent of $(u,g)$.
\end{lemma}

\begin{proof}
For an orthogonal matrix $O$, the process $\sigma\mapsto H_N(O\sigma)$ has the covariance of $H_N$, hence its law; so we may take $\sigma=\sqrt N\,e_N$ and $e_i$ the standard basis vectors. The coefficients of $H_N$ are Gaussian and $H_N$ is a polynomial, so all its derivatives are jointly Gaussian with $\E[\partial^\alpha H_N(x)\,\partial^\beta H_N(y)]=\partial_x^\alpha\partial_y^\beta C(x,y)$ for $C(x,y)=N^{1-p}\ip xy^p$. We differentiate $C$ and evaluate at $x=y=\sigma$, where $\ip\sigma\sigma=N$ and the coordinates $\sigma_i$, $i<N$, vanish; every term that keeps a factor $x_i$ or $y_i$ with $i<N$ vanishes. For $i,j,k,l<N$ this leaves
\begin{gather*}
\E H_N(\sigma)^2=N,\qquad\E[\partial_iH_N\,\partial_jH_N]=N^{1-p}\,p\,N^{p-1}\delta_{ij}=p\,\delta_{ij},\\
\E[\partial_{ij}H_N\,\partial_{kl}H_N]=N^{1-p}p(p-1)N^{p-2}(\delta_{ik}\delta_{jl}+\delta_{il}\delta_{jk})=\frac{p(p-1)}N(\delta_{ik}\delta_{jl}+\delta_{il}\delta_{jk}),
\end{gather*}
while $\E[H_N\,\partial_iH_N]$, $\E[H_N\,\partial_{ij}H_N]$ and $\E[\partial_kH_N\,\partial_{ij}H_N]$ are sums of such vanishing terms; for instance $\E[H_N(\sigma)\,\partial_iH_N(\sigma)]=N^{1-p}p\ip\sigma\sigma^{p-1}\sigma_i=0$ and $\E[H_N(\sigma)\,\partial_{ij}H_N(\sigma)]=N^{1-p}p(p-1)\ip\sigma\sigma^{p-2}\sigma_i\sigma_j=0$. Jointly Gaussian and uncorrelated, the three blocks are independent, and the covariance of $G$ is that of $\sqrt{p(p-1)}\,A$. By \eqref{eq:kacrice:hess} and Euler's identity, $\kacriceHess H_N(\sigma)=G-p\frac{H_N(\sigma)}N I=\sqrt{p(p-1)}\bigl(A-\frac{p}{\sqrt{p(p-1)}}uI\bigr)$, and $p/\sqrt{p(p-1)}=\sqrt{p/(p-1)}$.
\end{proof}

\begin{lemma}[one more eigenvalue]\label{lem:kacrice:extra}
There is a constant $c_N>0$ such that, for every $k\in\{0,\dots,N-1\}$ and every measurable $\varphi\ge0$,
\begin{equation}\label{eq:kacrice:extra}
\int_\R\varphi(t)\,e^{-Nt^2/4}\,\E\bigl[\abs{\det(A-tI)}\,\ind\{i(A-tI)=k\}\bigr]\dd t=\frac{c_N}N\,\E\varphi\bigl(\lambda_{(k)}\bigr),
\end{equation}
where $i(M)$ is the number of negative eigenvalues of $M$. Summing over $k$, $\int\varphi(t)e^{-Nt^2/4}\E\abs{\det(A-tI)}\dd t=c_N\int\varphi\rho_N$.
\end{lemma}

\begin{proof}
The plan is to write both sides with the joint density of the eigenvalues and to recognize $\abs{\det(A-t)}$ as the missing factors of a Vandermonde determinant. By Theorem~2.5.2 of \textcite{AndersonGuionnetZeitouni2010}, rescaled, the eigenvalues of $W$ listed in a uniformly random order have the density $q_N(x)=Z_N^{-1}\abs{\Delta(x)}\prod_{i=1}^Ne^{-Nx_i^2/4}$ on $\R^N$, with $\Delta(x)=\prod_{i<j}(x_j-x_i)$, and those of $A$ have the density $Z'^{-1}_{N-1}\abs{\Delta(y)}\prod_{i=1}^{N-1}e^{-Ny_i^2/4}$ on $\R^{N-1}$; both matrices have densities proportional to $e^{-N\tr(\cdot)^2/4}$. For eigenvalues $y_1,\dots,y_{N-1}$ of $A$, $\abs{\det(A-t)}=\prod_i\abs{y_i-t}$, and $i(A-t)=k$ means that exactly $k$ of them lie below $t$. Since $\abs{\Delta(y)}\prod_i\abs{y_i-t}=\abs{\Delta(y_1,\dots,y_{N-1},t)}$, Tonelli's theorem gives for the left side of \eqref{eq:kacrice:extra}
\[
\frac1{Z'_{N-1}}\int_{\R^N}\varphi(x_N)\,\ind\bigl\{\#\{i<N:x_i<x_N\}=k\bigr\}\,\abs{\Delta(x)}\prod_{i=1}^Ne^{-Nx_i^2/4}\dd x=\frac{Z_N}{Z'_{N-1}}\,\E\bigl[\varphi(\lambda'_N)\ind\{\lambda'_N=\lambda_{(k)}\}\bigr],
\]
where $(\lambda'_1,\dots,\lambda'_N)$ are the eigenvalues of $W$ in a uniformly random order; they are almost surely distinct. Given the set of eigenvalues, $\lambda'_N$ is each of them with probability $1/N$, so the expectation is $\E\varphi(\lambda_{(k)})/N$, and $c_N=Z_N/Z'_{N-1}$. For the sum, $\sum_k\E\varphi(\lambda_{(k)})=\E\sum_i\varphi(\lambda_i)=N\int\varphi\rho_N$; the density $\rho_N$ exists because $q_N$ does.
\end{proof}

\begin{lemma}[two spins]\label{lem:kacrice:two}
For $p=2$, almost surely $H_N$ has exactly $2N$ critical points on $S_N$, all nondegenerate.
\end{lemma}

\begin{proof}
For $p=2$, $H_N(\sigma)=\ip{\sigma}{B\sigma}$ with $B=(J+J^\top)/(2\sqrt N)$, which is $W/\sqrt2$ for a GOE matrix $W$. By \eqref{eq:kacrice:hess} the critical points are the $\sigma\in S_N$ with $B\sigma=\mu\sigma$ for some $\mu$, the eigenvectors of $B$ of norm $\sqrt N$. The eigenvalues of $W$ are almost surely distinct (\textcite{AndersonGuionnetZeitouni2010}, after Theorem~2.5.2), so there are exactly $2N$ of them, $\pm\sqrt Nv_i$, and at $\sqrt Nv_i$ the Hessian $2(B-\mu_iI)$ restricted to $v_i^\perp$ has the nonzero eigenvalues $2(\mu_j-\mu_i)$, $j\ne i$.
\end{proof}

\begin{theorem}[the exact mean count]\label{thm:kacrice:exact}
For all integers $p\ge2$, $N\ge2$, $k\in\{0,\dots,N-1\}$ and every Borel set $B\subset\R$,
\begin{gather}
\E\kacriceCrt_{N,k}(B)=2\sqrt{\tfrac2p}\,(p-1)^{N/2}\,\E\Bigl[e^{-N\frac{p-2}{4p}\lambda_{(k)}^2}\,\ind\bigl\{\lambda_{(k)}\in t(B)\bigr\}\Bigr],\label{eq:kacrice:exact}\\
\E\kacriceCrt_N(B)=2N\sqrt{\tfrac2p}\,(p-1)^{N/2}\int_{t(B)}e^{-N\frac{p-2}{4p}t^2}\rho_N(t)\dd t,\notag
\end{gather}
with $t(B)=\{t(u):u\in B\}=\sqrt{p/(p-1)}\,B$.
\end{theorem}

\begin{proof}
The plan is to apply the Kac--Rice formula, read off the integrand from \cref{lem:kacrice:regression}, convert the determinant with \cref{lem:kacrice:extra}, and find the constant from \cref{lem:kacrice:two}.

\emph{The formula on $S_N$.} \Cref{thm:toolkit:kacrice} applies to $f(y)=H_N(\sqrt Ny)$ on the unit sphere. Its hypotheses hold: in the chart $y\mapsto(y_1,\dots,y_{N-1})$ near $e_N$, the first and second partial derivatives at $e_N$ are the gradient and the Hessian of \cref{lem:kacrice:regression} (the term $\partial_Nf\cdot\partial_{ij}(1-\abs y^2)^{1/2}$ supplies $-pf\,\delta_{ij}$), whose joint law is nondegenerate by the lemma, so it stays nondegenerate near $e_N$ by continuity of the covariances, and finitely many rotations of this chart cover the sphere; the chart derivatives are polynomials in $y$ and $(1-\abs y^2)^{1/2}$ with Gaussian coefficients, so their covariances are smooth and hypothesis~(ii) holds with $\beta=1$. This is the verification of \textcite[proof of Theorem~2.1]{AuffingerBenArousCerny2013}. Under $\sigma=\sqrt Ny$ the gradient scales by $\sqrt N$, the Hessian by $N$ and the surface measure by $N^{(N-1)/2}$, and the factors $N^{N-1}$, $N^{-(N-1)/2}$ and $N^{-(N-1)/2}$ cancel: the formula holds on $S_N$ as it is written on the unit sphere.

\emph{The integrand.} By rotation invariance the integrand is the same at every point, and by \cref{lem:kacrice:regression} the gradient is independent of $(u,A)$, so conditioning on $g=0$ does not change their law and $p_\sigma(0)=(2\pi p)^{-(N-1)/2}$. With $\abs{S_N}=N^{(N-1)/2}\,2\pi^{N/2}/\Gamma(N/2)$ and $g_k(t)=\E[\abs{\det(A-tI)}\ind\{i(A-tI)=k\}]$,
\[
\E\kacriceCrt_{N,k}(B)=\abs{S_N}\,(2\pi p)^{-\frac{N-1}2}\bigl(p(p-1)\bigr)^{\frac{N-1}2}\int_B\sqrt{\tfrac N{2\pi}}\,e^{-Nu^2/2}g_k\bigl(t(u)\bigr)\dd u.
\]
\emph{The determinant.} Substitute $t=t(u)$: then $\dd u=\sqrt{(p-1)/p}\,\dd t$ and $\frac N2u^2=\frac N4t^2+N\frac{p-2}{4p}t^2$. \Cref{lem:kacrice:extra} with $\varphi(t)=e^{-N\frac{p-2}{4p}t^2}\ind\{t\in t(B)\}$ turns the integral into $\sqrt{(p-1)/p}\,\frac{c_N}N\E\varphi(\lambda_{(k)})$. The factors that depend on $p$ combine to $p^{-\frac{N-1}2}(p(p-1))^{\frac{N-1}2}\sqrt{(p-1)/p}=(p-1)^{N/2}/\sqrt p$, so $\E\kacriceCrt_{N,k}(B)=\gamma_N\,(p-1)^{N/2}p^{-1/2}\,\E\varphi(\lambda_{(k)})$ with a constant $\gamma_N$ that does not depend on $p$, $k$ or $B$.

\emph{The constant.} For $p=2$ the weight $\varphi$ is $\ind\{t\in t(B)\}$; with $B=\R$ and summed over $k$, the right side is $\gamma_N\cdot1\cdot2^{-1/2}\cdot N$, and the left side is $\E\kacriceCrt_N(\R)=2N$ by \cref{lem:kacrice:two}. Hence $\gamma_N=2\sqrt2$, which is \eqref{eq:kacrice:exact}; summing \eqref{eq:kacrice:exact} over $k$ with \cref{lem:kacrice:extra} gives the second formula.
\end{proof}

For $p=2$ the formula says $\E\kacriceCrt_N(B)=2N\int_{\sqrt2B}\rho_N$: the critical values of $\ip\sigma{B\sigma}$ are the eigenvalues of $W/\sqrt2$, each twice. Selberg's integral gives the same constant directly (\cref{ex:kacrice:mehta}).

\begin{theorem}[the annealed complexity]\label{thm:kacrice:annealed}
Let $p\ge2$ and $I_1(u)=J(2\abs u/\kacriceinf)$. For every $u\in\R$,
\[
\lim_{N\to\infty}\frac1N\log\E\kacriceCrt_N\bigl((-\infty,u)\bigr)=\Theta_p(u),\qquad\lim_{N\to\infty}\frac1N\log\E\kacriceCrt_{N,0}\bigl((-\infty,u)\bigr)=\Theta_{0,p}(u),
\]
where
\begin{gather*}
\Theta_p(u)=\begin{cases}\frac12\log(p-1)-\frac{p-2}{4(p-1)}u^2-I_1(u),&u\le-\kacriceinf,\\\frac12\log(p-1)-\frac{p-2}{4(p-1)}u^2,&-\kacriceinf\le u\le0,\\\frac12\log(p-1),&u\ge0,\end{cases}\\
\Theta_{0,p}(u)=\begin{cases}\frac12\log(p-1)-\frac{p-2}{4(p-1)}u^2-I_1(u),&u\le-\kacriceinf,\\\frac12\log(p-1)-\frac{p-2}p,&u\ge-\kacriceinf.\end{cases}
\end{gather*}
\end{theorem}

\begin{proof}
The plan is to evaluate the exponential rates of the two expectations in \eqref{eq:kacrice:exact} by bounding the weight $e^{-Nat^2}$ on the event that matters and by the large deviations \eqref{eq:kacrice:ldp}. Write $a=\frac{p-2}{4p}\ge0$ and $\tau=t(u)$, and note $a\tau^2=\frac{p-2}{4(p-1)}u^2$, $J(-\tau)=I_1(u)$ for $u\le-\kacriceinf$, and $4a=\frac{p-2}p$. Since $\frac1N\log\bigl(2\sqrt{2/p}\,N(p-1)^{N/2}\bigr)\to\frac12\log(p-1)$, \cref{thm:kacrice:exact} reduces the claims to the limits of $\frac1N\log Q_N$ and $\frac1N\log Q^0_N$, where
\[
Q_N=\int_{-\infty}^\tau e^{-Nat^2}\rho_N(t)\dd t=\E\Bigl[\frac1N\sum_ie^{-Na\lambda_i^2}\ind\{\lambda_i<\tau\}\Bigr],\qquad Q^0_N=\E\bigl[e^{-Na\lambda_{(0)}^2}\ind\{\lambda_{(0)}<\tau\}\bigr].
\]

\emph{All critical points, upper bounds.} If $\tau\ge0$, then $Q_N\le\E\mu_N(\R)=1$. If $\tau<0$, every term with $\lambda_i<\tau$ has $\lambda_i^2>\tau^2$, and at least one such term exists only if $\lambda_{(0)}<\tau$; hence $Q_N\le e^{-Na\tau^2}\Prob(\lambda_{(0)}<\tau)$. For $-2\le\tau<0$ we bound the probability by $1$; for $\tau<-2$ we use \eqref{eq:kacrice:ldp} for $-\lambda_{(0)}$, and $\limsup\frac1N\log Q_N\le-a\tau^2-J(-\tau)$.

\emph{All critical points, lower bounds.} If $-2<\tau\le0$, take $0<\delta<\tau+2$; on $(\tau-\delta,\tau)$ the weight is at least $e^{-Na(\tau-\delta)^2}$, so $Q_N\ge e^{-Na(\tau-\delta)^2}\E\mu_N((\tau-\delta,\tau))$, and $\E\mu_N((\tau-\delta,\tau))\to\sigma_{\mathrm{sc}}((\tau-\delta,\tau))>0$ by the semicircle law (\cref{thm:toolkit:semicircle}) and bounded convergence. Hence $\liminf\frac1N\log Q_N\ge-a(\tau-\delta)^2$ for every small $\delta$. If $\tau>0$, the interval $(-\delta,0)$ gives $\liminf\ge-a\delta^2$ in the same way. If $\tau\le-2$, then on the event $\{\tau-\delta<\lambda_{(0)}<\tau\}$ the sum defining $Q_N$ is at least $\frac1Ne^{-Na(\tau-\delta)^2}$, so $Q_N\ge\frac1Ne^{-Na(\tau-\delta)^2}\Prob(-\tau<-\lambda_{(0)}<-\tau+\delta)$, and \eqref{eq:kacrice:ldp} gives $\liminf\frac1N\log Q_N\ge-a(\tau-\delta)^2-J(-\tau)$. Letting $\delta\to0$ matches the upper bounds in all cases.

\emph{Minima.} For $\tau\le-2$ the same two estimates apply verbatim to $Q^0_N$: $Q^0_N\le e^{-Na\tau^2}\Prob(\lambda_{(0)}<\tau)$ and $Q^0_N\ge e^{-Na(\tau-\delta)^2}\Prob(\tau-\delta<\lambda_{(0)}<\tau)$. For $\tau>-2$ and $0<\delta<\min(2,\tau+2)$, split according to whether $\lambda_{(0)}\le-2+\delta$: on that event $\lambda_{(0)}^2\ge(2-\delta)^2$, so
\[
Q^0_N\le e^{-Na(2-\delta)^2}+\Prob\bigl(\lambda_{(0)}>-2+\delta\bigr),
\]
and the last probability is $\Prob(\lambda_{\max}(W)<2-\delta)$, whose logarithm divided by $N$ tends to $-\infty$ by \eqref{eq:kacrice:ldp}; hence $\limsup\frac1N\log Q^0_N\le-a(2-\delta)^2$. In the other direction $(-2-\delta,-2)\subset(-\infty,\tau)$, so $Q^0_N\ge e^{-Na(2+\delta)^2}\Prob(2<-\lambda_{(0)}<2+\delta)$ and $\liminf\frac1N\log Q^0_N\ge-a(2+\delta)^2-J(2)=-a(2+\delta)^2$. Letting $\delta\to0$ gives the limit $-4a$.
\end{proof}

Both functions are continuous and nondecreasing. On $u\le0$, $\Theta_p(u)$ is also the rate of the number of critical points with energy in a small window around $u$, so it is the physicist's $\Sigma(u)$ of \eqref{eq:kacrice:sigma}, and below the threshold $\Theta_{0,p}=\Theta_p$: on the exponential scale, almost every critical point below $-\kacriceinf$ is a minimum. The proof shows more than the limits: the weight $e^{-Nat^2}$ favors critical values near $0$, the entropy of the eigenvalue favors the bulk, and the threshold is where the eigenvalue variable leaves the bulk.

\begin{corollary}[the threshold and the ground state]\label{cor:kacrice:groundstate}
Let $p\ge3$.
\begin{enumerate}
\item For every $u>-\kacriceinf$, $\frac1N\log\E\kacriceCrt_{N,0}((u,\infty))\to-\infty$.
\item The function $\Theta_{0,p}$ is strictly increasing on $(-\infty,-\kacriceinf]$ and has a unique zero $-E_0$ there, with $E_0>\kacriceinf$. Almost surely $\liminf_N\kacriceGS_N\ge-E_0$, where $\kacriceGS_N=\frac1N\min_{S_N}H_N$.
\end{enumerate}
\end{corollary}

\begin{proof}
(i) By \eqref{eq:kacrice:exact} with $B=(u,\infty)$, $\E\kacriceCrt_{N,0}((u,\infty))\le2\sqrt{2/p}\,(p-1)^{N/2}\Prob(\lambda_{(0)}>t(u))$, and $t(u)>-2$, so the probability is $\Prob(\lambda_{\max}(W)<-t(u))$ with $-t(u)<2$, and \eqref{eq:kacrice:ldp} applies.

(ii) On $u<-\kacriceinf$ the term $-\frac{p-2}{4(p-1)}u^2$ is strictly increasing for $p\ge3$, and $I_1(u)=J(\abs{t(u)})$ is decreasing in $u$ because $J$ increases on $[2,\infty)$; so $\Theta_{0,p}$ is strictly increasing there and tends to $-\infty$ as $u\to-\infty$. At $u=-\kacriceinf$ it equals $\frac12\log(p-1)-\frac{p-2}p$, which is positive for $p\ge3$: it is $0.0132$ at $p=3$, and its derivative in $p$, $\frac1{2(p-1)}-\frac2{p^2}=\frac{(p-2)^2}{2p^2(p-1)}$, is nonnegative. So the zero $-E_0$ exists, is unique, and lies below $-\kacriceinf$. Fix $\eps>0$. A global minimizer of $H_N$ on $S_N$ is a critical point of index $0$, so by Markov's inequality and \cref{thm:kacrice:annealed},
\[
\Prob\bigl(\kacriceGS_N<-E_0-\eps\bigr)\le\E\kacriceCrt_{N,0}\bigl((-\infty,-E_0-\eps)\bigr)=e^{N(\Theta_{0,p}(-E_0-\eps)+o(1))},
\]
and $\Theta_{0,p}(-E_0-\eps)<0$. The bounds are summable, and the Borel--Cantelli lemma gives $\kacriceGS_N\ge-E_0-\eps$ for all large $N$, almost surely.
\end{proof}

\subsection{What the general theorems add}\label{sec:kacrice:general}

\begin{theorem}[fixed index; \textcite{AuffingerBenArousCerny2013}, Theorems~2.5 and~2.14--2.15]\label{thm:kacrice:index}
For $p\ge2$ and every fixed $k\ge0$, $\frac1N\log\E\kacriceCrt_{N,k}((-\infty,u))\to\Theta_{k,p}(u)$, where $\Theta_{k,p}$ is $\Theta_{0,p}$ with $(k+1)I_1(u)$ in place of $I_1(u)$. With $E_k$ the zero of $\Theta_{k,p}(-\,\cdot)$, $E_0>E_1>\dots\to\kacriceinf$; for every $\eps>0$ the probability that a critical point of index $k$ lies above $-\kacriceinf+\eps$ is at most $e^{-cN^2}$, and the probability that one of index at least $k$ lies below $-E_k-\eps$ decays exponentially.
\end{theorem}

\begin{proofsketch}
The first statement is the proof of \cref{thm:kacrice:annealed} with $\lambda_{(k)}$ in place of $\lambda_{(0)}$: by Theorem~A.1 of \textcite{AuffingerBenArousCerny2013}, which extends \textcite{BenArousDemboGuionnet2001}, the $(k+1)$-th smallest eigenvalue obeys a large deviation principle with rate $(k+1)J(-\,\cdot)$, because $k+1$ eigenvalues must leave the bulk. The $e^{-cN^2}$ bound replaces the speed-$N$ estimate of \cref{cor:kacrice:groundstate}(i) by the large deviation principle of the empirical measure at speed $N^2$ \parencite{BenArousGuionnet1997b}; the last statement is Markov's inequality, as in the corollary. What is omitted is the large deviation principle for $\lambda_{(k)}$.
\end{proofsketch}

\begin{theorem}[the typical count; \textcite{Subag2017a}, Theorems~1 and~3, Corollary~2]\label{thm:kacrice:subag}
Let $p\ge3$. For every $u\in(-E_0,-\kacriceinf)$, $\E[\kacriceCrt_N((-\infty,u))^2]/(\E\kacriceCrt_N((-\infty,u)))^2\to1$, so $\kacriceCrt_N((-\infty,u))/\E\kacriceCrt_N((-\infty,u))\to1$ in $L^2$ and in probability. For every $u>-E_0$ the two moments agree on the exponential scale: $\frac1N\log\E[\kacriceCrt_N((-\infty,u))^2]\to2\Theta_p(u)$.
\end{theorem}

\begin{proofsketch}
A two-point Kac--Rice formula writes $\E[\kacriceCrt_N^2]$ as an integral over pairs of points, organized by their overlap $r$; given that both gradients vanish, the two Hessians are correlated Gaussian matrices, whose determinants Subag controls. Below the threshold the contribution of pairs with overlap away from $0$ is exponentially smaller than $(\E\kacriceCrt_N)^2$, and the contribution of nearly orthogonal pairs matches it up to a factor $1+o(1)$. Omitted: the analysis of the two-point Hessian and of the overlaps near $\pm1$; Section~3 of \textcite{Subag2017a} outlines the proofs, which occupy Sections~7 and~8.
\end{proofsketch}

Together with \cref{cor:kacrice:groundstate}, the second moment gives $\kacriceGS_N\to-E_0$ for every $p\ge3$ \parencite[Section~1 and Appendix~IV]{Subag2017a}; \textcite[Theorem~2.12]{AuffingerBenArousCerny2013} had proved it for even $p\ge4$ through the Parisi formula for spherical models \parencite{Talagrand2006b}, and zero-temperature Parisi formulas \parencite{ChenSen2017,JagannathTobasco2017b} give the same value. Three routes to $E_0(3)$ agree to twelve digits: the zero of $\Theta_{0,3}$, the zero-temperature 1RSB infimum of \cref{sec:kacrice:replicas}, and the ground state of \textcite[Theorem~1.2]{JagannathLopattoMiolane2020} at $\lambda=0$. Beyond the mean, \textcite{SubagZeitouni2017} show that the critical values near the minimum form, after centering, a Poisson process with exponential intensity, so that the minimum fluctuates like a Gumbel variable, and \textcite{Subag2017b} shows that at low temperature the Gibbs measure lives on thin bands around the deepest minima. Above the threshold, \textcite{SubagZeitouni2021} prove concentration at every energy above $-E_0$ for $p\ge32$. For mixed models the annealed count can mislead: \textcite[Corollary~4.1]{AuffingerBenArous2013} exhibit models in which the mean number of critical points below an energy grows exponentially while the probability that one exists tends to zero. Kac--Rice also bounds the TAP complexity of an inference problem: \textcite[Theorem~1.1]{FanMeiMontanari2021} bound the expected number of critical points of the TAP free energy of $\Z_2$ synchronization (\cref{ch:amp}) and deduce that its low-energy local minimizers are close to the posterior mean when the signal is strong. On the algorithmic side, \textcite{HuangSellke2025} prove that for even $p\ge4$ an algorithm with Lipschitz dependence on the couplings reaches energy below $-\kacriceinf-\eps$ only with exponentially small probability (their Theorem~5 and Proposition~2.2 in arXiv v2): for such algorithms the threshold is the limit.\sidenote{For the pure model with even $p\ge4$, the value $\mathsf{ALG}$ of \textcite[Proposition~2.2]{HuangSellke2025} is $\int_0^1\xi''(q)^{1/2}\dd q$ with $\xi(q)=q^p$, which is $2\sqrt{(p-1)/p}=\kacriceinf$; the algorithm of \textcite{Subag2021} reaches it.}

\section{The proof, continued: the spiked tensor}\label{sec:kacrice:spiked}

The spike changes the Kac--Rice integrand by deterministic shifts, which the next proposition computes exactly at every $N$. It is Lemma~2 of \textcite{BenArousMeiMontanariNica2019} in our normalization.

\begin{proposition}[the spiked count at finite $N$]\label{prop:kacrice:spiked}
Let $p\ge3$, $\lambda\ge0$, and let $\kacriceCrt_{N,\lambda}(M,B)$ count the critical points of $H_{N,\lambda}$ with $m(\sigma)\in M$ and $H_{N,\lambda}(\sigma)/N\in B$, for Borel sets $M\subset(-1,1)$ and $B\subset\R$. Then
\[
\E\kacriceCrt_{N,\lambda}(M,B)=K_N\int_M(1-m^2)^{\frac{N-3}2}e^{-\frac N2p\lambda^2m^{2p-2}(1-m^2)}\int_Be^{-\frac N2(u+\lambda m^p)^2}\,\E\bigl\lvert\det\bigl(A-t(u)I-\theta(m)ee^\top\bigr)\bigr\rvert\dd u\dd m,
\]
with $\theta(m)=\lambda\sqrt{p(p-1)}\,m^{p-2}(1-m^2)$, a fixed unit vector $e\in\R^{N-1}$, and the constant $K_N=\abs{S_N}\frac{\Gamma(N/2)}{\sqrt\pi\,\Gamma((N-1)/2)}\bigl(\frac{p-1}{2\pi}\bigr)^{\frac{N-1}2}\sqrt{\frac N{2\pi}}$. The number of local minima has the same formula with the indicator of $A-t(u)I-\theta(m)ee^\top\succ0$ inside the expectation.
\end{proposition}

\begin{proof}
We use the Kac--Rice formula for $H_{N,\lambda}$, a centered Gaussian field plus the smooth function $-\lambda Nm(\sigma)^p$, as \textcite[Lemmas~1--2]{BenArousMeiMontanariNica2019} do; the rest is the regression of \cref{lem:kacrice:regression} shifted by the spike. At a point $\sigma$ with $m(\sigma)=m$, the spike $F_s(\sigma)=-\lambda N^{1-p}\ip\sigma{x^*}^p$ has gradient $-\lambda pm^{p-1}x^*$ and Hessian $-\lambda p(p-1)m^{p-2}x^*x^{*\top}/N$, and $\ip\sigma{\nabla F_s}=-\lambda pNm^p$. The tangential part of $x^*$ has norm $\sqrt{N(1-m^2)}$; let $e$ be its direction. By \eqref{eq:kacrice:hess}, along the sphere $F_s$ has gradient $-\lambda pm^{p-1}\sqrt{N(1-m^2)}\,e$ and Hessian $-\lambda p(p-1)m^{p-2}(1-m^2)ee^\top+\lambda pm^pI$.

The gradient of $H_{N,\lambda}$ vanishes when the gradient $g\sim\Normal(0,pI_{N-1})$ of $H_N$ equals $\lambda pm^{p-1}\sqrt{N(1-m^2)}\,e$, whose density is $(2\pi p)^{-(N-1)/2}e^{-Np\lambda^2m^{2p-2}(1-m^2)/2}$; since $g$ is independent of the value and the second derivatives of $H_N$, the conditioning leaves their law unchanged. The energy $u=H_{N,\lambda}/N$ equals $H_N/N-\lambda m^p$, whose density is $\sqrt{N/(2\pi)}\,e^{-N(u+\lambda m^p)^2/2}$. The Hessian is
\[
\sqrt{p(p-1)}\,A-p\bigl(u+\lambda m^p\bigr)I-\lambda p(p-1)m^{p-2}(1-m^2)ee^\top+\lambda pm^pI=\sqrt{p(p-1)}\bigl(A-t(u)I-\theta(m)ee^\top\bigr),
\]
and by the rotation invariance of $A$ we may fix $e$. Finally, the surface measure of $S_N$ pushed forward by $m$ is $\abs{S_N}\,\frac{\Gamma(N/2)}{\sqrt\pi\,\Gamma((N-1)/2)}(1-m^2)^{(N-3)/2}\dd m$, the law of one coordinate of a uniform point of the unit sphere; collecting the constants as in \cref{thm:kacrice:exact} gives $K_N$.
\end{proof}

At $m=0$ and $p\ge3$ the spike disappears from the integrand, which is the physicist's observation that the equator does not see the signal. The asymptotics need the exponential rate of $\E\abs{\det(A-t-\theta ee^\top)}$ and of the probability that this matrix is positive definite, where the outlier of \cref{ch:spiked} enters.

\begin{theorem}[the landscape of the spiked tensor; \textcite{BenArousMeiMontanariNica2019}, Theorems~1--2 and Propositions~1--2, in our normalization]\label{thm:kacrice:bammn}
Let $p\ge3$ and $\lambda\ge0$ be fixed. For Borel sets $M\subset[-1,1]$ and $B\subset\R$,
\[
\sup_{m\in M^\circ,u\in B^\circ}\Sigma_\lambda(m,u)\le\liminf_{N\to\infty}\frac1N\log\E\kacriceCrt_{N,\lambda}(M,B)\le\limsup_{N\to\infty}\frac1N\log\E\kacriceCrt_{N,\lambda}(M,B)\le\sup_{m\in\bar M,u\in\bar B}\Sigma_\lambda(m,u),
\]
with $\Sigma_\lambda$ as in \eqref{eq:kacrice:sigmalambda}, and the same holds for local minima with $\Sigma_\lambda(m,u)-L(\theta(m),-t(u))$, where $L(\theta,s)=+\infty$ for $s<2$; $L(\theta,s)=\frac12J(s)-\frac12J(b)-\frac\theta2(s-b)+\frac18(s^2-b^2)$ for $\theta>1$ and $2\le s<b=\theta+\frac1\theta$; and $L=0$ otherwise. The function $S(m)=\sup_u\Sigma_\lambda(m,u)$ satisfies, for $0\le m<m_c=((p-2)/(\lambda\sqrt{p(p-1)}))^{1/p}$,
\[
S(m)=\tfrac12\log(p-1)+\tfrac12\log(1-m^2)-\tfrac p2\lambda^2m^{2p-2}(1-m^2)+\frac{p\lambda^2}{2(p-2)}m^{2p};
\]
and $S(m)\le0$ on $[m_c,1)$, where $S$ vanishes exactly at the largest root $m^*$ of $m^{2p-4}(1-m^2)=1/(p\lambda^2)$, which exists if and only if $\lambda\ge\lambda_c(p)$ of \eqref{eq:kacrice:lambdac}. On the exponential scale, critical points with overlap in $[m_c,1)$ exist in mean only near $m^*$, and only for $\lambda\ge\lambda_c(p)$.
\end{theorem}

\begin{proofsketch}
By \cref{prop:kacrice:spiked} the count is a Laplace integral in $(m,u)$; the factors $(1-m^2)^{(N-3)/2}$, the gradient density and the energy density give the first four terms of $\Sigma_\lambda$, and $\abs{S_N}$ with the constants gives $\frac12(1+\log(p-1))$. A rank-one perturbation changes the determinant only by the factor $\abs{1-\theta\ip e{(A-t)^{-1}e}}$, which is not exponential, so $\frac1N\log\E\abs{\det(A-t-\theta ee^\top)}\to\Phi(t)$; for minima, the outlier must be pushed above $t$, at the cost $L$ given by the large deviations of the largest eigenvalue of a rank-one perturbation of the GOE \parencite{Maida2007}. Omitted: the uniformity in $(m,u)$ needed for Laplace's method, and the control of $\E\abs{\det}$ near the outlier; see Section~4 of \textcite{BenArousMeiMontanariNica2019}. The formulas for $S$ are elementary calculus (their Appendix~B).
\end{proofsketch}

The equatorial band is now a theorem about means: $S(0)=\frac12\log(p-1)$, and the complexity of minima at $m=0$ is $\Theta_{0,p}(-\kacriceinf)$, for every $\lambda$. The estimator itself is understood exactly.

\begin{theorem}[the maximum likelihood estimator; \textcite{JagannathLopattoMiolane2020}, Theorem~1.2]\label{thm:kacrice:jlm}
Let $p\ge3$, $f_\lambda(q)=\lambda^2q^p+\log(1-q)+q$ on $[0,1)$, $\kacriceML=\sup\{\lambda\ge0:\sup_qf_\lambda(q)\le0\}$, and $q^*(\lambda)=0$ for $\lambda<\kacriceML$ and $q^*(\lambda)=\arg\max f_\lambda$ for $\lambda>\kacriceML$. If $\hat\sigma$ minimizes $H_{N,\lambda}$ on $S_N$, then almost surely $\abs{m(\hat\sigma)}\to\sqrt{q^*(\lambda)}$ for every $\lambda\ne\kacriceML$.
\end{theorem}

The stationarity of $f_\lambda$ is $q^{p-2}(1-q)=1/(p\lambda^2)$: above $\kacriceML$, the overlap of the estimator is the location $m^*$ of the good critical point of \cref{thm:kacrice:bammn}. At $\kacriceML$ it jumps from $0$ to $m^*(\kacriceML)$, a first-order transition; for $p=3$, $\kacriceML=1.20656$ (with $q=0.645007$) and $m^*(\kacriceML)=0.80312$. Both thresholds agree with the replica predictions of \textcite{RosEtAl2019} quoted in the physics box: $\lambda_c(3)=1.15470$ and $\kacriceML(3)=1.20656$, against their $r_c=\sqrt6$ and $r_{\mathrm{1ST}}\approx2.56$ (that is, $1.2068$). For $\lambda_c<\lambda<\kacriceML$ the good minimum exists but is not the deepest. And between the statistical threshold, of order one, and the algorithmic one lies the equatorial band: gradient flow and Langevin dynamics started uniformly on the hemisphere $m>0$ recover the spike in time of order one when $\lambda=N^\alpha$ with $\alpha>\frac{p-2}2$ \parencite[Theorem~2.1]{BenArousGheissariJagannath2020}; for $\alpha<\frac{p-2}2$, Langevin dynamics at positive temperature started from a Gibbs-type measure close to the uniform one has exponentially small probability, in a power of $N$, of reaching overlap $2N^{-1/2+\eps}$ within a time $e^{cN^\eps}$ (their Theorem~2.10). Their proof suggests that the obstruction is the weak curvature of the spike near the equator rather than the count of minima, a point we return to in \cref{sec:kacrice:fails}.

\section{Measuring it}\label{sec:kacrice:measuring}

Four scripts test the formulas. \texttt{13-kacrice-complexity.py} evaluates every curve and constant of this chapter: $\kacriceinf$, the $E_k$, the agreement of the three routes to $E_0$, the misprint check of \cref{sec:kacrice:notes}, $\lambda_c$, $\kacriceML$ and $m^*$, and the curves of \cref{fig:kacrice:complexity} and \cref{fig:kacrice:spiked}(a), in about a minute.

\paragraph{The exact formula at small $N$.} \texttt{13-kacrice-count.py} finds every critical point of the $3$-spin on $S_N$ for $N=3,\dots,7$: Newton's method on the equations $\nabla H_N=\mu\sigma$, $\norm\sigma^2=N$, started from the endpoints of a gradient-descent sweep and then from random points, in batches until three batches in a row find no new critical value. Every instance satisfies the Morse identity $\sum_k(-1)^k\kacriceCrt_{N,k}(\R)=1+(-1)^{N-1}$; for odd $N$ this certifies that no pair $\pm\sigma$ of critical points was missed. The average counts are compared with \eqref{eq:kacrice:exact}, whose GOE expectations are computed from $400000$ matrices. \Cref{fig:kacrice:count}(a) shows the agreement for all critical points and for minima: every interval contains the formula. Index by index the agreement is the same; at $N=7$ the formula gives $4.83$, $10.69$, $15.72$, $17.71$ critical points of indices $0$ to $3$, and the count $4.84$, $10.60$, $15.58$, $17.63$. Without the descent sweep, random-start Newton missed about $6\%$ of the minima at $N=7$, an error that the Morse identity catches only in part: it cannot see a minimum and a saddle of index $1$ missed together with their antipodes. The counts are small. A generic cubic form has at most $2(2^N-1)$ critical points on the sphere, twice the number of eigenvalues of a generic symmetric tensor of order $3$ \parencite{CartwrightSturmfels2013}: $14$ at $N=3$, which some instances attain, and $254$ at $N=7$, where the mean is $80$.

\begin{figure}[t]
\begin{fullwidth}
\centering
\begin{tikzpicture}
\begin{groupplot}[group style={group size=2 by 1,horizontal sep=18mm},rplot,width=82mm,height=52mm]
\nextgroupplot[xlabel={dimension $N$},ylabel={mean number},ymode=log,xmin=2.6,xmax=7.4,ymin=2,ymax=150,
  xtick={3,4,5,6,7},ytick={2,5,10,20,50,100},yticklabels={2,5,10,20,50,100},
  title={(a) the exact formula at small $N$},legend pos=north west,legend style={font=\scriptsize}]
\addplot[proofpred,mark=none] table[x=N,y=formula_all] {figures/data/13-kacrice-count.dat};
\addplot[proofpred,mark=none,forget plot] table[x=N,y=formula_min] {figures/data/13-kacrice-count.dat};
\addplot[measured,ebars] table[x=N,y=count_all,y error plus=half_all,y error minus=half_all] {figures/data/13-kacrice-count.dat};
\addplot[measured3,ebars] table[x=N,y=count_min,y error plus=half_min,y error minus=half_min] {figures/data/13-kacrice-count.dat};
\legend{\cref{thm:kacrice:exact},all critical points,local minima}
\nextgroupplot[xlabel={dimension $N$},ylabel={energy per spin},xmode=log,xmin=13,xmax=330,ymin=-1.68,ymax=-1.28,
  xtick={16,32,64,128,256},xticklabels={16,32,64,128,256},ytick={-1.65,-1.6,-1.5,-1.4,-1.3},
  title={(b) where the minima are, and where descent stops},legend pos=north east,legend style={font=\scriptsize}]
\addplot[physpred] coordinates {(13,-1.6329932) (330,-1.6329932)};
\addplot[guide] coordinates {(13,-1.6569984) (330,-1.6569984)};
\addplot[proofpred,mark=*,mark size=1.2pt,mark options={solid,fill=proof}] table[x=N,y=umean] {figures/data/13-kacrice-minima.dat};
\addplot[measured,ebars] table[x=N,y=energy,y error plus=half,y error minus=half] {figures/data/13-kacrice-gd.dat};
\legend{threshold $-E_\infty$,ground state $-E_0$,mean minimum (exact formula),gradient descent}
\end{groupplot}
\end{tikzpicture}
\end{fullwidth}
\caption{The pure spherical $3$-spin at finite $N$. (a) Mean numbers of critical points and of local minima: \cref{thm:kacrice:exact} evaluated with $400000$ GOE matrices (lapis dashed; Monte Carlo error below $0.02\%$) against exhaustive counts on $200$, $150$, $120$, $100$ and $80$ random Hamiltonians at $N=3,\dots,7$ (seeds $1300+N$), with $95\%$ Student-$t$ intervals. (b) The mean energy of a local minimum in the expected count, from \cref{thm:kacrice:exact} with $40000$ samples of the smallest GOE eigenvalue per $N$ (lapis; Monte Carlo errors below $0.002$; seed $1313$), and the energy at which projected gradient descent from uniform random starts stops (points: $95\%$ Student-$t$ intervals over $10$, $8$, $6$ and $5$ Hamiltonians with $8$ starts each, at $N=25$, $50$, $75$, $100$; seeds $130+N$), against the threshold of \cref{thm:kacrice:annealed} (carnelian) and the ground state (dashed). Scripts: \texttt{13-kacrice-count.py}, \texttt{13-kacrice-minima.py}, \texttt{13-kacrice-descent.py}.}\label{fig:kacrice:count}
\end{figure}

\paragraph{Where the minima are at finite $N$.} The exponent $\Theta_{0,3}(-\kacriceinf)=0.0132$ is small: at $N=128$ the formula gives only $455$ local minima on average. At such sizes the asymptotic picture does not yet hold, and the exact formula says where the minima are. Applied to $g(u)=u$, \eqref{eq:kacrice:exact} gives the mean energy of a local minimum in the expected count, $\E\sum_{\text{minima}}H_N/N$ divided by $\E\kacriceCrt_{N,0}(\R)$, as a ratio of two expectations over the smallest eigenvalue of a GOE matrix. \texttt{13-kacrice-minima.py} samples that eigenvalue with the tridiagonal model of \textcite{DumitriuEdelman2002} and a Sturm-sequence bisection. The mean energy of a minimum is $-1.319$, $-1.415$, $-1.482$, $-1.526$ and $-1.557$ at $N=16$, $32$, $64$, $128$ and $256$, far above $-\kacriceinf=-1.633$. The share of the expected count below $-\kacriceinf$ is $1.8\%$ at $N=16$ and $0.02\%$ at $N=256$, and the gap to the threshold decreases like $1.2/\sqrt N$: the product of gap and $\sqrt N$ stays between $1.21$ and $1.26$. A heuristic explains the rate. For a minimum at $t=-2+x$ inside the bulk, the weight $e^{-Nat^2}$ gains $e^{4aNx}$, while the probability that the smallest eigenvalue of $A$ exceeds $-2+x$ costs about $e^{-cN^2x^3}$; the balance is at $x\asymp N^{-1/2}$.

\paragraph{Where descent stops.} \texttt{13-kacrice-descent.py} runs projected gradient descent with step $0.15$ until the gradient falls below $10^{-5}\sqrt N$. The endpoints have mean energy $-1.449$, $-1.511$, $-1.537$ and $-1.558$ at $N=25$, $50$, $75$ and $100$ (\cref{fig:kacrice:count}(b)). Descent stops at deeper minima than the average one of the exact formula at the same $N$, by $0.05$ to $0.07$, which is natural for minima with large basins; but it follows the same slow approach to $-\kacriceinf$ and stays far above the ground state. Single runs scatter widely at small $N$ (from $-1.735$ to $-1.181$ at $N=25$, where the ground state itself fluctuates far from $-E_0$); at $N=100$ the lowest of the $40$ endpoints is $-1.643$. At the endpoints the Hessian is the semicircle of \cref{lem:kacrice:regression} centered at the energy (\cref{fig:kacrice:hessian}), except at the lower edge, which a minimum must keep nonnegative. At $N=100$ the mean eigenvalue is $1.932$, against $-t(\bar u)=1.908$ at the mean energy $\bar u=-1.558$, and the largest eigenvalues (median $3.84$) sit near the edge $2-t(\bar u)=3.91$; but the semicircle at that energy reaches down to $-0.09$, while every endpoint has a positive smallest eigenvalue, with median $0.079$ and range $[0.021,0.166]$. Descent stops at minima that are as close to marginal as the size allows.

\begin{figure}[t]
\centering
\begin{tikzpicture}
\begin{axis}[rplot,xlabel={eigenvalue of $\kacriceHess H_N/\sqrt6$},ylabel={density},xmin=-0.5,xmax=4.5,ymin=0,ymax=0.36,
  legend pos=north east,legend style={font=\scriptsize}]
\addplot[ybar interval,fill=stone!25,draw=stone!60,forget plot] table[x=x,y=density] {figures/data/13-kacrice-hessian.dat};
\addplot[physpred] table[x=x,y=energy] {figures/data/13-kacrice-hessian-sc.dat};
\addplot[proofpred] table[x=x,y=threshold] {figures/data/13-kacrice-hessian-sc.dat};
\addplot[guide,forget plot] coordinates {(0,0) (0,0.36)};
\legend{semicircle at the endpoints' energy,semicircle at $-E_\infty$}
\end{axis}
\end{tikzpicture}
\caption{The Hessian at the $40$ local minima found by gradient descent at $N=100$ (histogram of the $99\times40$ eigenvalues of the Hessian along the sphere, divided by $\sqrt{p(p-1)}=\sqrt6$), against the semicircle on $[-2-t,2-t]$ that \cref{lem:kacrice:regression} predicts at the mean energy $u$ of the endpoints, $t=2u/E_\infty$ (carnelian), and the semicircle of the threshold, on $[0,4]$ (lapis dashed). Same runs as \cref{fig:kacrice:count}(b). Script: \texttt{13-kacrice-descent.py}.}\label{fig:kacrice:hessian}
\end{figure}

\paragraph{The spiked tensor.} \texttt{13-kacrice-spiked.py} runs the same descent on $H_{N,\lambda}$ at $N=80$, with the step reduced to $0.3/\lambda$ for $\lambda>2$, where the spike stiffens the landscape (\cref{fig:kacrice:spiked}(b)). From warm starts at overlap $0.5$, descent lands on the good minimum: for $\lambda=1.2$, $1.5$, $2$, $3$ and $6$ the final overlap is $0.810$, $0.918$, $0.959$, $0.983$ and $0.996$, against $m^*=0.798$, $0.905$, $0.953$, $0.981$ and $0.995$, with intervals below $0.01$; the excess, largest near $\lambda_c$, is a finite-size effect. Below $\lambda_c$ the warm starts stop at minima of intermediate overlap, between $0.22$ and $0.60$ at $\lambda=0.9$ and between $0.33$ and $0.75$ at $\lambda=1.1$, beyond the edge of the band where the annealed complexity of minima is positive ($m\approx0.24$ and $0.34$). There that complexity is negative but tiny, about $-0.001$ at $\lambda=1.1$, so that $e^{80\Sigma}\approx0.9$: at this size such minima are not rare. At $\lambda=0.6$ all ten warm starts end at $\abs m\le0.27$, near the band. From random starts the overlap begins at about $N^{-1/2}=0.11$, and descent reaches the spike in $0$ of $10$ runs at $\lambda\le0.9$, in between $1$ and $4$ of $10$ for $1.1\le\lambda\le3$, and in $8$ and $7$ of $10$ at $\lambda=4$ and $6$; at this size $N^{(p-2)/2}=9$, so the algorithmic scale is close.

\section{The status board}\label{sec:kacrice:status}

For the pure spherical $p$-spin, the landscape statements of physics are theorems for the mean counts at every energy and every fixed index, and for the typical count below the threshold. The ground state is known for every $p\ge3$. What is open is dynamical: that gradient descent ends at the threshold is a conjecture, proved for low-temperature Langevin dynamics on time scales of order one (\cref{ch:dmft}). For the spiked tensor, the mean landscape and the estimator are theorems, the typical landscape is a physical prediction, and the algorithmic threshold is proved in one direction.

\begin{statusboard}
\statusrow{Mean number of critical points below energy $u$: all, minima, fixed index $k$; no fixed-index critical points above $-\kacriceinf$}{\proved}{\textcite[Theorems~2.5, 2.8, 2.14, 2.15]{AuffingerBenArousCerny2013}; \cref{thm:kacrice:annealed} (all critical points and minima)}
\statusrow{Typical number $=$ mean number (pure model)}{\provedunder}{below $-\kacriceinf$ for $p\ge3$ \parencite[Theorem~1]{Subag2017a}; at all energies above $-E_0$ for $p\ge32$ \parencite{SubagZeitouni2021}}
\statusrow{Typical number $=$ mean number (mixed models)}{\knownfalse}{full mixtures \parencite[Corollary~4.1]{AuffingerBenArous2013}}
\statusrow{Ground state $-E_0$, zero of the complexity of minima}{\proved}{lower bound: \cref{cor:kacrice:groundstate}; equality: \textcite{AuffingerBenArousCerny2013} (even $p$), \textcite{Subag2017a} (all $p\ge3$)}
\statusrow{Lowest critical point of index $k$ at $-E_k$}{\conjectured}{\textcite[Remark~2.16]{AuffingerBenArousCerny2013}; the layers are proved as probability bounds}
\statusrow{Gradient descent from a random start ends at $-\kacriceinf$}{\conjectured}{\textcite{CugliandoloKurchan1993}; numerics (\cref{fig:kacrice:count}); the Langevin analogue as $\beta\to\infty$ is proved \parencite{Sellke2024}}
\statusrow{Spiked tensor: mean landscape, $\lambda_c(p)$, equatorial band}{\proved}{\textcite[Theorems~1--2, Propositions~1--2]{BenArousMeiMontanariNica2019}}
\statusrow{Spiked tensor: typical landscape, thresholds $\lambda_c$, $\kacriceML$}{\conjectured}{replicated Kac--Rice at the RS level \parencite{RosEtAl2019}; both thresholds match proved mean or estimator results}
\statusrow{Estimator threshold $\kacriceML$ and overlap $\sqrt{q^*(\lambda)}$}{\proved}{\textcite[Theorem~1.2]{JagannathLopattoMiolane2020}}
\statusrow{Descent from a random start needs $\lambda\gtrsim N^{(p-2)/2}$}{\provedunder}{sufficiency \parencite[Theorem~2.1]{BenArousGheissariJagannath2020}; necessity for Langevin from a Gibbs-type start (their Theorem~2.10), conjectured from the uniform start}
\end{statusboard}

\section{Where the translation fails}\label{sec:kacrice:fails}

\paragraph{What was assumed.} The model is pure and the disorder Gaussian. The exact formula rests on Euler's identity $\ip{\sigma}{\nabla H_N}=pH_N$, which makes the Hessian at energy $u$ a GOE matrix shifted by a deterministic amount; for mixed models the shift is random given the energy, and the formulas of \textcite{AuffingerBenArous2013} are more involved. It rests also on the joint density of GOE eigenvalues, which non-Gaussian couplings do not have. Two 2026 preprints take up the annealed count for non-Gaussian couplings: \textcite{ChenLuSen2026} prove that it keeps its Gaussian limit when the couplings mix a strictly subgaussian law with a nonzero Gaussian component (they write that one step of their proof \enquote{was suggested by ChatGPT (OpenAI)}), and \textcite{Kim2026} constructs a symmetric law of bounded support, matching the first $2p$ Gaussian moments, for which the mean count in a compact energy window grows at a strictly larger exponential rate. Every limit is taken with $\lambda$, $p$ and the index fixed as $N\to\infty$; the algorithmically relevant regime $\lambda=N^\alpha$ lies outside the theorems on the landscape. The theorems count critical points in mean; the physics asks for typical counts, and the gap is closed only for the pure model below the threshold, or at all energies for $p\ge32$.

\paragraph{Where the physics is wrong, or not yet right.} The annealed count is wrong as a typical count for full mixtures \parencite{AuffingerBenArous2013}. For mixed models the threshold picture of part~(iii) fails as a description of dynamics: numerical integration of the dynamical equations shows gradient descent from equilibrated initial conditions going below the threshold energy of the dominant marginal minima \parencite{FolenaFranzRicciTersenghi2020}. For the spiked tensor, \textcite{RosEtAl2019} report that the thermodynamic replica method gives partly incorrect predictions for the landscape, which their replicated Kac--Rice computation corrects; and \textcite{BenArousGheissariJagannath2020} argue that the failure of descent for $\lambda\ll N^{(p-2)/2}$ comes from the weakness of the signal's curvature near the equator, not from the count of minima, so that the physicist's explanation through the rugged landscape is not the mechanism of their proof. Finally, at accessible sizes the asymptotic landscape is a poor guide: at $N\le256$ the exact formula places almost all minima above $-\kacriceinf$ (\cref{fig:kacrice:count}).

\paragraph{Open problems.} As far as we know, the following are open as of October 2026. (1) The typical complexity of the spiked tensor, in particular near the isolated minimum and in the band near $\lambda_c$: \textcite{RosEtAl2019} expect a rigorous version of their replica-symmetric computation, and \textcite{BenArousMeiMontanariNica2019} leave typical counts open. (2) Concentration of the number of critical points above the threshold for $3\le p\le31$, which the second moment of \textcite{SubagZeitouni2021} reaches only for $p\ge32$. (3) The energy of the lowest critical point of index $k\ge1$, conjectured to tend to $-E_k$ \parencite[Remark~2.16]{AuffingerBenArousCerny2013}. (4) That gradient flow, not only low-temperature Langevin dynamics, ends at $-\kacriceinf$ on the pure model; \textcite[Remark~1.5]{Sellke2024} explains why his method does not reach it. (5) Refutation below $N^{(p-2)/2}$ for descent started from the uniform measure, which \textcite[Section~2.4]{BenArousGheissariJagannath2020} believe to hold.

\begin{exercises}
\item \consolidation (The covariances.) Complete the proof of \cref{lem:kacrice:regression}: compute $\partial_{x_i}\partial_{x_j}\partial_{y_k}C$ and $\partial_{x_i}\partial_{x_j}\partial_{y_k}\partial_{y_l}C$ for $C(x,y)=N^{1-p}\ip xy^p$ and check that at $x=y=\sqrt Ne_N$ only $\frac{p(p-1)}N(\delta_{ik}\delta_{jl}+\delta_{il}\delta_{jk})$ survives.\label{ex:kacrice:cov}

\item \consolidation (Two spins.) Show directly that for $p=2$ the formula of \cref{thm:kacrice:exact} reads $\E\kacriceCrt_{N,k}(B)=2\Prob(\lambda_{(k)}\in\sqrt2B)$, and explain it with \cref{lem:kacrice:two}.\label{ex:kacrice:two}

\item \consolidation (Counting with signs.) (a) Show that $\E\det(tI-A)=N^{-(N-1)/2}\mathrm{He}_{N-1}(t\sqrt N)$, with $\mathrm{He}_n$ the probabilists' Hermite polynomial: in the expansion of the determinant only fixed points and transpositions survive. (b) With the proof of \cref{thm:kacrice:exact}, show $\sum_k(-1)^k\E\kacriceCrt_{N,k}(\R)=1+(-1)^{N-1}$ for every $p\ge2$, the Euler characteristic of the sphere. \hint{$\E\mathrm{He}_{2m}(cZ)=(2m-1)!!\,(c^2-1)^m$ for $Z\sim\Normal(0,1)$.}\label{ex:kacrice:euler}

\item \consolidation (The constants.) Check the numbers of \cref{comp:kacrice:potential}: $\Phi(0)=-\frac12$, $\Sigma(-\kacriceinf)=\frac12\log(p-1)-\frac{p-2}p$, and $E_0(3)=1.656998$, $E_0(4)=1.794085$.\label{ex:kacrice:numbers}

\item \byhand (Selberg's constant.) Mehta's integral $\int_{\R^n}\abs{\Delta(x)}\prod_ie^{-x_i^2/2}\dd x=(2\pi)^{n/2}\prod_{j=1}^n\Gamma(1+\frac j2)/\Gamma(\frac32)$ gives $Z_N$ and $Z'_{N-1}$ in \cref{lem:kacrice:extra}. Show $c_N=(2/N)^{N/2}\,2\sqrt2\,\Gamma(1+\frac N2)$ and recover $\gamma_N=2\sqrt2$ in the proof of \cref{thm:kacrice:exact} without \cref{lem:kacrice:two}.\label{ex:kacrice:mehta}

\item \byhand (The equator.) For $p\ge3$ show that the gradient and the Hessian of $-\lambda Nm(\sigma)^p$ along the sphere vanish at $m=0$, and deduce from \cref{thm:kacrice:bammn} that $S(0)=\frac12\log(p-1)$ and that the complexity of minima at $m=0$ is $\frac12\log(p-1)-\frac{p-2}p$ for every $\lambda$.\label{ex:kacrice:equator}

\item \byhand (The good critical point.) (a) Show that $m^{2p-4}(1-m^2)$ is maximal at $m^2=\frac{p-2}{p-1}$ and deduce \eqref{eq:kacrice:lambdac}; check $\lambda_c(3)=2/\sqrt3$, $\lambda_c(4)=\sqrt{27}/4$. (b) Show that at $\lambda=\lambda_c$ the root $m^*$ equals $m_c$ of \cref{thm:kacrice:bammn}. (c) For $p=3$ show $m^*(\lambda)^2=\frac12(1+\sqrt{1-4/(3\lambda^2)})$.\label{ex:kacrice:mstar}

\item \byhand (The estimator's threshold.) Show that $\kacriceML$ of \cref{thm:kacrice:jlm} is $(pq^{p-2}(1-q))^{-1/2}$ at the root $q\in(0,1)$ of $\frac{q^2}{p(1-q)}+q+\log(1-q)=0$, and compute $\kacriceML(3)=1.20656$, $\kacriceML(4)=1.40567$.\label{ex:kacrice:jlm}

\item \translate (No minima above the threshold, at speed $N^2$.) Using the large deviation principle of the empirical measure at speed $N^2$ (\cref{rem:toolkit:goeldp}), prove that $\Prob(\lambda_{(0)}(W)>-2+\delta)\le e^{-c_\delta N^2}$ and deduce that the probability of a local minimum above $-\kacriceinf+\eps$ is at most $e^{-cN^2}$. Which sentence of the physicist's argument does this replace?\label{ex:kacrice:n2}

\item \translate (Quenched equals annealed.) Let $X_N\ge0$ with $\E X_N^2/(\E X_N)^2\to1$. Show $X_N/\E X_N\to1$ in probability, and that $\Prob(X_N>0)\ge(\E X_N)^2/\E X_N^2$. Combine with \cref{thm:kacrice:subag} and the concentration of $\kacriceGS_N$ (\cref{thm:toolkit:lipschitz}) to prove $\kacriceGS_N\to-E_0$. Which replica assumption does each step replace?\label{ex:kacrice:second}

\item \translate (TAP states.) From \eqref{eq:kacrice:tap}, show that the stationary points in $q$ at energy $h=H_N(\sigma)/N$ satisfy, with $z=(1-q)q^{p/2-1}$, $\beta z=\frac1{\sqrt2(p-1)}(-h\pm\sqrt{h^2-\kacriceinf^2})$, so that they exist only for $h\le-\kacriceinf$ \parencite[Section~6]{AuffingerBenArousCerny2013}. What does this say about TAP states above the threshold?\label{ex:kacrice:tap}

\item \measurement (The approach to the threshold.) In \texttt{13-kacrice-minima.py}, add $N=512$ and $1024$ with more samples, and test whether $(\bar u_N+\kacriceinf)\sqrt N$ stays near $1.2$. Why does the effective sample size of the importance weights shrink with $N$, and what would a better sampler of the smallest eigenvalue condition on?\label{ex:kacrice:approach}

\item \measurement (Escaping the equator.) In \texttt{13-kacrice-spiked.py}, find for $N=40$, $80$ and $160$ the $\lambda$ at which half of the random starts reach $m>\frac12$, and compare its growth with $N^{(p-2)/2}=\sqrt N$.\label{ex:kacrice:escape}

\item \research (The typical spiked landscape.) Compute the second moment of the number of critical points of $H_{N,\lambda}$ near the equator for $\lambda<\lambda_c$, and decide whether the band of minima there is typical, as \textcite{RosEtAl2019} predict.\label{ex:kacrice:research}
\end{exercises}

\notesandrefs\label{sec:kacrice:notes}
Counting critical points of random functions began with the real roots of random polynomials \parencite{Kac1943} and the level crossings of random noise \parencite{Rice1944}; \textcite{AdlerTaylor2007} and \textcite{AzaisWschebor2009} give the modern theory. In spin glasses, \textcite{BrayMoore1980} counted TAP solutions of the Sherrington--Kirkpatrick model, and for the spherical $p$-spin \textcite{KurchanParisiVirasoro1993}, \textcite{CrisantiSommers1995}, \textcite{CavagnaGiardinaParisi1998} and \textcite{CrisantiLeuzziRizzo2003} computed the complexity, while \textcite{CugliandoloKurchan1993} found the threshold in the dynamics. \textcite{Fyodorov2004} reduced the mean count of an isotropic landscape to the mean absolute determinant of a GOE-type matrix and solved it at finite $N$ for a model of a manifold in a random environment, and \textcite{BrayDean2007} related index and energy for isotropic fields; \textcite{CastellaniCavagna2005} give a pedagogical account of the physics.

The rigorous theory is due to \textcite{AuffingerBenArousCerny2013}, whose identity between the absolute determinant and one more eigenvalue is the heart of \cref{thm:kacrice:exact}, and was extended to mixed models by \textcite{AuffingerBenArous2013}. The determination of the constant through $p=2$, and the restriction to the cases that need only the largest eigenvalue, are the book's choices. Subag's second moment \parencite{Subag2017a} made the ground state accessible without the Parisi formula and started a fine geometric theory of the pure model \parencite{Subag2017b,SubagZeitouni2017,SubagZeitouni2021}. For the spiked tensor, \textcite{MontanariRichard2014} introduced the model and its algorithms, \textcite{BenArousMeiMontanariNica2019} computed the mean landscape, \textcite{RosEtAl2019} the typical one by physics, \textcite{JagannathLopattoMiolane2020} the estimator, and \textcite{BenArousGheissariJagannath2020} the thresholds of descent.

The measurements were made for this book: the exhaustive counts at small $N$, the finite-$N$ energy profile of the minima computed from the exact formula (\cref{fig:kacrice:count}(b)), its comparison with gradient descent, and the descent on the spiked tensor, whose warm-start overlaps \cref{thm:kacrice:bammn,thm:kacrice:jlm} predict. The tridiagonal sampler is that of \textcite{DumitriuEdelman2002}.

\paragraph{Translating the literature.} \Textcite{AuffingerBenArousCerny2013} and \textcite{Subag2017a} use our $H_N$, $\kacriceCrt_{N,k}$, $\Theta_p$, $\Theta_{k,p}$, $E_k$ and $\kacriceinf$, but normalize the GOE with $\E M_{ij}^2=(1+\delta_{ij})/(2N)$, so their $\lambda_k^N$ is our $\lambda_{(k)}/\sqrt2$ and their eigenvalue set $\sqrt{p/(2(p-1))}B$ is our $t(B)/\sqrt2$. Their Remark~2.7 states the Crisanti--Sommers complexity as $\Theta_{0,p}(2^{-1/2}E)$; it should read $\Theta_{0,p}(\sqrt2E)$, as their Lemma~6.3 has it. We checked this to ten digits (\texttt{13-kacrice-complexity.py}): with $2^{-1/2}E$ the right side is constant below the threshold. Physicists \parencite{CrisantiSommers1995,CastellaniCavagna2005,CugliandoloKurchan1993} use energies $E=u/\sqrt2$, so the threshold of the $3$-spin is $-1.154701$ and its ground state $-1.171675$, and temperatures $\sqrt2$ times smaller than ours. \Textcite{BenArousMeiMontanariNica2019} maximize $f=-H_{N,\lambda}/(\sqrt2N)$ on the unit sphere, so their $\lambda$ is our $\lambda/\sqrt2$, their energy $x$ is our $-u/\sqrt2$, their $\lambda_c(3)=\sqrt{2/3}$ is our $2/\sqrt3$, and their local maxima are our minima. \Textcite{JagannathLopattoMiolane2020} and \textcite{BenArousGheissariJagannath2020} use our $\lambda$; the former write $\lambda_s$ for our $\lambda_c$ and $\lambda_c$ for our $\kacriceML$. \Textcite{RosEtAl2019} use physics energies and the spike $rNm^k/k$, so their $r$ is $p\lambda/\sqrt2$. \Textcite{SubagZeitouni2021} write $E_\infty$ for the ground state, our $E_0$; \textcite{Sellke2024} and \textcite{HuangSellke2025} maximize $H_N$.

\printchapterbib
\end{refsection}
